\documentclass[11pt,a4paper]{article}
\usepackage[utf8]{inputenc}
% lmodern must precede fontenc: without a scalable Type1 font the T1 encoding
% falls back to bitmap fonts, and microtype's expansion then aborts the run
% with "auto expansion is only possible with scalable fonts".
\usepackage{lmodern}
\usepackage[T1]{fontenc}
\usepackage[margin=1in]{geometry}
\usepackage{amsmath,amssymb}
\usepackage{graphicx}
\usepackage{microtype}
\usepackage[hidelinks,breaklinks]{hyperref}
\usepackage{url}
\urlstyle{same}

% Generated by code/tools/md_to_tex.py from the SurveyForge Markdown output.
% Citation numbers match the source .md exactly.

\title{Retrieval-Augmented Generation for Large Language Models: A Comprehensive Survey of Foundations, Architectures, and Emerging Frontiers}
\author{Generated by SurveyForge}
\date{}

\begin{document}
\maketitle
\tableofcontents
\newpage



\section{Introduction}


The remarkable fluency and breadth of large language models (LLMs) often mask a fundamental architectural limitation: their knowledge is entirely encoded in static, parametric weights. This gives rise to persistent challenges---most notably factual hallucination, where models generate plausible but incorrect or unsupported assertions \cite{arxiv:2005.11401v4,arxiv:2401.01313v3}. The issue is compounded by temporal knowledge cutoffs; once trained, an LLM cannot account for events or discoveries that occur after its data collection, creating a critical need for updatable world knowledge \cite{arxiv:2212.10511v4,arxiv:2310.07521v3}. Furthermore, purely parametric models often lack the depth of domain-specific expertise required in specialized fields, and the prohibitive computational cost and risk of catastrophic forgetting during frequent retraining preclude continuous refreshment \cite{arxiv:2312.05934v3,arxiv:2402.01364v2}. These limitations have motivated a paradigm shift toward semi-parametric architectures that decouple knowledge storage in a non-parametric, easily modifiable external corpus from the parametric language model, giving rise to Retrieval-Augmented Generation (RAG).

RAG fundamentally redefines the generation process by conditioning an LLM on relevant documents retrieved from an external knowledge base at inference time \cite{arxiv:2005.11401v4}. In contrast to fine-tuning, which bakes new information into the model weights and often fails on low-frequency facts, RAG provides instance-level access to external memory without altering the underlying model \cite{arxiv:2402.11457v1}. It also differs from model editing, which surgically modifies a small number of parameters to update specific facts and can still suffer from generalization and consistency issues, and from continual pre-training, which incrementally updates all parameters at considerable expense \cite{arxiv:2402.01364v2,arxiv:2402.04624v1}. This semi-parametric design enables RAG systems to remain transparent, adaptable, and---crucially---attributable, offering a practical path toward more trustworthy AI.

The historical progression toward modern RAG reflects a convergence of information retrieval and language generation. Early work on retrieval-assisted text generation laid the conceptual groundwork \cite{arxiv:2202.01110v2}, but the inflection came with the introduction of differentiable retrieval mechanisms and dense vector representations. The pioneering RAG model \cite{arxiv:2005.11401v4} and RETRO \cite{arxiv:2112.04426v3} demonstrated that combining pre-trained retrievers with parametric generators could yield strong performance on knowledge-intensive benchmarks. The development of powerful dense retrievers built on pretrained transformers further fueled adoption, enabling semantic matching at scale \cite{arxiv:2211.14876v1}. A subsequent wave of modular, black-box augmentation techniques such as REPLUG \cite{arxiv:2301.12652v4} and self-reflective strategies like Self-RAG \cite{arxiv:2310.11511v1} shifted the emphasis toward flexible, adaptive retrieval-generation loops. Concurrently, end-to-end training methods that relax independence assumptions, exemplified by Stochastic RAG \cite{arxiv:2405.02816v1}, have further strengthened the synergy between retriever and generator.

This survey provides a holistic, multi-dimensional taxonomy of the rapidly expanding RAG landscape. We systematically analyze the retrieval backbone---spanning sparse, dense, and hybrid models---the design of knowledge sources, and the integration architectures that fuse retrieved evidence into generation, ranging from input-level concatenation to deep latent-space fusion. We then examine training and optimization paradigms, including joint end-to-end learning, instruction tuning, and alignment with human preferences. Evaluation methodologies, from component-level metrics to integrated faithfulness benchmarks, are critically reviewed, alongside the robustness, security, and privacy implications of external data access. The survey further maps domain-specific adaptations---in healthcare, law, code generation, and conversational AI---and identifies open challenges and emerging frontiers. By synthesizing the foundational techniques and charting future directions, we aim to equip researchers and practitioners with a rigorous, unified framework for advancing retrieval-augmented LLMs.


\section{Foundations and Taxonomy of Retrieval-Augmented Generation}



\subsection{Formal Problem Definition and the Generic Retrieval-Augmented Generation Pipeline}


This subsection establishes a rigorous mathematical and architectural foundation for retrieval-augmented generation (RAG), decomposing the paradigm into a probabilistic objective, its constituent components, and the end-to-end information flow that distinguishes it from purely parametric generation. At its core, RAG formalizes text generation as a process conditioned on both the input query and externally retrieved evidence, thereby mitigating factual hallucination and knowledge cutoffs inherent in standalone large language models (LLMs) \cite{arxiv:2005.11401v4,arxiv:2112.04426v3}.

The joint retrieval-generation process is canonically modeled as a probabilistic graphical model. Given an input sequence \( x \), the objective is to generate a target output \( y \) by marginalizing over a latent variable \( z \), which represents retrieved evidence from an external knowledge store \( \mathcal{Z} \):  
\[
P(y|x) = \sum_{z \in \mathcal{Z}} P_{\text{gen}}(y|x, z) P_{\text{ret}}(z|x)
\]  
Here, \( P_{\text{ret}}(z|x) \) embodies the retrieval component, scoring the relevance of evidence \( z \) to the input, while \( P_{\text{gen}}(y|x, z) \) is the generation component, a parametric model that synthesizes the final output by conditioning on both the original input and the retrieved evidence. This formulation, originating in the original RAG framework \cite{arxiv:2005.11401v4}, frames RAG as a semi-parametric model: the retrieval prior \( P_{\text{ret}} \) provides a non-parametric interface to a dynamic external corpus, while \( P_{\text{gen}} \) remains a parametric sequence-to-sequence or decoder-only language model. Alternative formulations treat the retrieval process as a deterministic, hard-selection mechanism where a top-\( k \) set of documents is passed to the generator, a simplification that facilitates black-box integration with commercial LLMs \cite{arxiv:2301.12652v4} but sacrifices the probabilistic marginalization over evidence.

The architecture is instantiated through three core component definitions. The \textbf{retriever} is a function \( \mathcal{R}: x \to \mathcal{Z} \) that maps an input query to a set of relevant documents or passages from the external store. This mapping can be realized through sparse lexical models (e.g., BM25), dense dual-encoder architectures (e.g., DPR) \cite{arxiv:2211.14876v1}, or hybrid ensembles that combine complementary signals. The \textbf{external knowledge store} \( \mathcal{D} \) is a non-parametric corpus, typically a collection of unstructured text passages indexed for efficient similarity search. Its construction, indexing strategy, and maintenance profoundly influence system performance, as the retriever must efficiently search and retrieve high-quality evidence at inference time. The \textbf{generator} is a parametric LLM that consumes a context formed by the concatenation or fusion of \( x \) and the retrieved set \( z \), and then generates the output \( y \). A critical design distinction exists between black-box generator access, where only the input-output interface is available, and white-box access, where internal representations and gradients can be manipulated \cite{arxiv:2301.12652v4,arxiv:2310.11511v1}. Black-box approaches, such as those that prepend retrieved passages to the prompt, rely entirely on the model's in-context learning capabilities, while white-box approaches permit deeper architectural fusion, such as cross-attention conditioning on the retrieved chunks \cite{arxiv:2112.04426v3} or joint end-to-end training \cite{arxiv:2208.03299v3}.

The end-to-end information flow through a generic RAG pipeline is characterized by four sequential stages, each introducing distinct computational and qualitative considerations. \textbf{Indexing} is an offline stage that pre-processes and encodes the knowledge store into a searchable representation, constructing vector indices or inverted indices to enable efficient retrieval. \textbf{Retrieval} is the online stage where the input query is encoded and matched against the index to fetch the top-\( k \) most relevant items according to a similarity scoring function. \textbf{Augmentation} then injects this retrieved content into the generation context, which can be accomplished through simple input-level concatenation, cross-attention within the decoder layers \cite{arxiv:2112.04426v3}, or latent-space fusion. Finally, \textbf{synthesis} is the generation stage where the LLM autoregressively decodes the output, attending to both the original query and the augmented evidence. This pipeline, however, is susceptible to cascading failures: a retrieval failure where \( P_{\text{ret}}(z|x) \) assigns high probability to irrelevant documents can introduce distracting noise, while an integration failure where \( P_{\text{gen}} \) fails to correctly condition on faithful evidence can lead to hallucination, despite correct retrieval \cite{arxiv:2310.01558v1,arxiv:2401.05856v1}. This motivates the evolution from a static, single-shot pipeline toward more dynamic and adaptive instantiations, such as iterative retrieval-generation loops \cite{arxiv:2305.06983v2} and self-reflective systems that critique and refine both the retrieved evidence and the generated output \cite{arxiv:2310.11511v1}, a trajectory that fundamentally redefines the flow of information and control in modern RAG systems.


\subsection{Core Component Anatomy and Design Decisions}


The design of a retrieval-augmented generation (RAG) system is profoundly shaped by architectural decisions concerning its constituent components. While the formal pipeline defines the process, the anatomy of each stage---specifically, what constitutes a unit of knowledge, what types of knowledge are accessible, and when retrieval is invoked---determines the system's fidelity, latency, and overall effectiveness. This subsection dissects these core design dimensions: retrieval granularity, knowledge source characteristics, and the temporal dynamics of retrieval relative to generation.

A foundational decision is the unit of retrieval, or granularity, which fundamentally dictates context density and the signal-to-noise ratio for the generator. Passage-level retrieval, as employed in foundational RAG models \cite{arxiv:2005.11401v4}, provides broad contextual blocks that can support open-domain question answering but often introduce superfluous information that can distract the generator. In contrast, sentence-level granularity, pursued in some iterative systems, offers more precise evidence pinpointing but may lack the broader coherence necessary for complex reasoning. Entity-centric and token-level retrieval represent the extreme end of this spectrum, with methods like xRAG reinterpreting dense document embeddings as a distinct modality to achieve extreme context compression down to a single token, thereby eliminating the need for explicit text expansion \cite{arxiv:2405.13792v1}. Emerging hierarchical strategies seek to balance these trade-offs, employing multi-level indices that first retrieve coarse-grained passages and then refine the context to specific sentences or facts in a cascading manner, optimizing both recall and relevance.

Beyond text snippets, the nature of the external knowledge source profoundly dictates indexing strategies and the encoding of evidence for the generator. While unstructured text corpora like Wikipedia remain the most common repository \cite{arxiv:2005.11401v4,arxiv:2208.03299v3}, systems increasingly integrate heterogeneous knowledge. Retrieval from structured databases and Knowledge Graphs (KGs) introduces structured relational information, enabling precise querying of entities and their connections. GraphRAG systems formalize this by incorporating Graph-Based Indexing and Graph-Enhanced Generation to capture multi-hop relational knowledge that is difficult to infer from flat text \cite{arxiv:2408.08921v2}. Similarly, the extension to multimodal knowledge---encompassing images, tables, and code---is critical for tasks beyond text-only question answering. Models like MuRAG and RA-CM3 access external multimodal memory, retrieving both text and images to generate enriched, grounded outputs \cite{arxiv:2210.02928v2,arxiv:2211.12561v2}. The design challenge shifts to building unified retrieval interfaces and fusion modules that can seamlessly reason over this heterogeneous evidence, a capability for which unified frameworks like UniMS-RAG are being developed \cite{arxiv:2401.13256v1}.

The third critical dimension is the timing of retrieval relative to the generation process, which distinguishes static, single-shot augmentation from dynamic, interactive paradigms. The dominant approach is pre-generation, one-shot retrieval, where a query is used once to fetch a set of documents that are then prepended to the generator's input \cite{arxiv:2301.12652v4}. This simple, modular design is computationally efficient but fails for long-form or multi-hop tasks where evolving information needs arise mid-generation. To address this, iterative retrieval interleaves multiple retrieval steps with generation, allowing the model to formulate new queries based on its partially generated output. Frameworks like Iter-RetGen \cite{arxiv:2305.15294v2} and FLARE \cite{arxiv:2305.06983v2} exemplify this synergy, where a prediction of future content or a low-confidence token triggers a new retrieval cycle, progressively gathering evidence. The most sophisticated systems employ adaptive retrieval, where the model itself decides dynamically whether to consult its parametric knowledge or seek external evidence. This is often gated by confidence estimations or special reflection tokens, as seen in Self-RAG and metacognitive approaches like MetaRAG, which enable the system to monitor, evaluate, and plan its retrieval strategy introspectively \cite{arxiv:2402.11626v1}. This adaptive paradigm transforms the system from a static pipeline into an autonomous agent that actively manages its cognitive resources, trading off latency for enhanced factual grounding on an as-needed basis. 

Together, the choices of retrieval granularity, knowledge source type, and retrieval timing establish the baseline trade-off space for a RAG system's efficiency, factual precision, and ability to handle complex tasks. However, these component-level decisions represent only one facet of the broader RAG design landscape. The following section builds upon these dimensions to introduce complementary axes---integration depth and architectural coupling---that, together with retrieval timing, form a comprehensive taxonomy for understanding and navigating modern RAG architectures.


\subsection{Multi-Dimensional Taxonomy of Retrieval-Generation Paradigms}


The design space of retrieval-augmented generation systems can be systematically organized along three fundamental axes: retrieval timing, integration depth, and architectural coupling. This multi-dimensional taxonomy provides a comparative framework for analyzing the inherent trade-offs---accuracy versus latency, flexibility versus synergy, and transparency versus performance---that shape modern RAG architectures.

\textbf{Retrieval timing} captures when and how frequently external knowledge is fetched relative to the generation process. The simplest, most widespread paradigm is single-shot retrieval, where a query is issued once before generation begins and the resulting passages are fed to the language model as a static context. While computationally efficient, this approach struggles with complex, multi-step reasoning where later evidence needs depend on earlier generation outputs. To address this, iterative multi-hop retrieval interleaves retrieval and generation: partial outputs drive new queries, progressively accumulating relevant information. Systems such as IRCoT \cite{arxiv:2212.10509v2} weave retrieval steps into chain-of-thought reasoning, while Iter-RetGen \cite{arxiv:2305.15294v2} reprompts the model with all previously retrieved knowledge to produce an improved answer. Moving beyond fixed schedules, adaptive retrieval decides dynamically \emph{whether} to retrieve based on model confidence or detected knowledge gaps. Techniques like CRAG \cite{arxiv:2401.15884v2} employ a lightweight evaluator to assess retrieval quality and optionally extend search to the web, while other works \cite{arxiv:2402.13492v3} adapt retrieval frequency according to the popularity of queried entities and relations. These three timing modes trade off latency, reasoning depth, and robustness: single-shot is fast but brittle, iterative gains accuracy at the cost of multiple retrieval rounds, and adaptive attempts to balance both by retrieving only when necessary, though it relies on reliable confidence estimation.

\textbf{Integration depth} describes how retrieved evidence is fused with the generative model's internal representations. At the shallowest end, input-level concatenation prepends retrieved passages directly into the model's prompt. This method is model-agnostic, requires no architectural changes, and is the backbone of most deployed RAG systems, but it suffers from positional biases and finite context-window limits. Deeper integration is achieved through cross-attention conditioning, as epitomized by FiD (Fusion-in-Decoder) \cite{arxiv:2209.14290v1}, where the decoder attends to all encoded passages independently, alleviating concatenation bottlenecks and enabling more thorough evidence grounding. Further along the continuum, latent-space fusion injects retrieval representations directly into intermediate model layers, bypassing text entirely. xRAG \cite{arxiv:2405.13792v1} reinterprets dense document embeddings as a distinct modality and fuses them via a trainable bridge, achieving extreme compression while preserving retrieval semantics. At the deepest level, virtual token modulation prepends learnable soft embeddings to the prompt, steering the generation process without modifying the base LLM's weights. The choice of depth directly impacts the system's flexibility: shallow fusion works with any black-box generator, whereas deeper fusion requires white-box access but yields tighter coupling and often higher factual accuracy.

\textbf{Architectural coupling} concerns the degree of co-adaption and information flow between the retriever and the generator. In tightly joint systems, both components are trained end-to-end, enabling gradient signals from generation loss to update retrieval parameters. Foundational models such as REALM and RAG, as well as more recent unified architectures like CorpusLM \cite{arxiv:2402.01176v2}, exemplify this paradigm, achieving strong synergy but at the expense of high training cost and limited component replaceability. In contrast, modular designs treat the retriever and generator as independent, potentially black-box modules. This decoupling simplifies deployment, allows off-the-shelf components to be exchanged, and is widely used in commercial RAG pipelines; however, it risks a mismatch between the retriever's relevance score and the generator's actual information need. A middle ground is represented by systems like Re2G \cite{arxiv:2207.06300v1}, which combine a separate retriever, reranker, and generator but train them jointly through knowledge distillation. Within this axis, the distinction between white-box and black-box generators further refines the design possibilities: white-box access enables architectural fusion (cross-attention, adapters), whereas black-box only permits input-level augmentation.

These three axes are not independent; they interact to define the overall system behavior. For instance, a tightly coupled system with iterative retrieval and latent-space fusion can achieve state-of-the-art factual accuracy, whereas a modular single-shot pipeline with input-level augmentation prioritizes simplicity and scalability. The taxonomy underscores that no single paradigm universally dominates---the optimal configuration is dictated by task requirements, latency budgets, model accessibility, and the tolerance for training complexity. Emerging work increasingly explores hybrid designs that dynamically select integration depth or retrieval timing based on query difficulty, pointing toward a future of RAG systems that adaptively navigate this multi-dimensional design space.


\subsection{Comparative Positioning Among Knowledge Infusion Techniques}


Building on the preceding taxonomy, where architectural coupling and integration depth chart a spectrum from strictly modular to deeply fused systems, RAG's semi-parametric design occupies the extreme modular position. It maintains an explicit separation between non-parametric external knowledge and parametric model memory, enabling decoupled, independently updatable access to information without altering the language model's weights. This architectural commitment fundamentally distinguishes RAG from other knowledge infusion strategies---fine-tuning, knowledge editing, and in-context learning---that rely on parametric recall or prompt-encapsulated knowledge to varying degrees. Below we position RAG relative to these paradigms and highlight how it can be synergistically combined with them to build more robust systems.

\textbf{Fine-tuning and continual pre-training} embed new knowledge directly into the weights of a large language model (LLM) through gradient updates on domain-specific corpora. While this can yield highly efficient, context-na\"{i}ve inference and strong domain adaptation, it is a static and expensive process: any change in the knowledge base necessitates retraining or additional fine-tuning, risking catastrophic forgetting of previously learned information. RAG, by contrast, operates on-the-fly, retrieving knowledge from a continuously updatable external store without modifying model parameters \cite{arxiv:2312.10997v5}. This makes RAG especially advantageous for rapidly evolving data (e.g., news, financial reports) or tail entities whose infrequent occurrence renders fine-tuning impractical. Nonetheless, fine-tuning can be more effective when the target distribution is stable and retrieval noise is high; it guarantees that all relevant background knowledge is implicitly encoded without the latency and potential irrelevance of an external search. In practice, the two techniques are complementary: a fine-tuned model can serve as a strong generator within an RAG pipeline, reducing sensitivity to retrieval errors while still benefiting from external evidence.

\textbf{Knowledge editing} techniques (e.g., ROME, MEMIT) offer a middle ground by surgically modifying a handful of model weights to update specific factual associations. These methods provide persistence and precision comparable to fine-tuning but at a fraction of the computational cost and with localized impact that aims to avoid general forgetting. However, they still alter the parametric knowledge of the model, which can lead to unintended side effects and scaling difficulties when many edits are required. RAG bypasses this risk entirely by never modifying the generator; it supplies corrective information as context, and the model's original world knowledge remains intact. The boundary between the two is blurring with the emergence of \emph{editable external memory} structures that combine a stored knowledge base with selective overwriting mechanisms, thereby merging the persistence of parametric editing with the flexibility of retrieval. While such approaches are not yet mainstream, they underscore a future direction where a retrieval corpus itself becomes the target of lightweight, targeted updates.

\textbf{In-context learning and long-context models} offer another alternative by packing all necessary knowledge into the prompt, either as few-shot demonstrations or by presenting the entire relevant corpus within an extended context window. RAG fundamentally differs in its ability to scale to corpora that vastly exceed the model's maximum context length, as it retrieves only a handful of relevant passages rather than processing the full database. Long-context LLMs can, in principle, read entire long documents, avoiding retrieval latency and potential ranking errors, but they suffer from quadratic attention complexity and the ``lost-in-the-middle'' phenomenon where information in the middle of a long input is poorly utilized. RAG sidesteps these issues by using retrieval as a pre-filter, feeding the generator only the most pertinent snippets. Moreover, RAG can be applied to black-box LLMs without access to internal representations, as demonstrated by REPLUG \cite{arxiv:2301.12652v4}, which simply prepends retrieved documents to the input. Recent work such as LongRAG \cite{arxiv:2406.15319v3} highlights a convergence: using a long-context LLM as the generator in a RAG loop, where larger retrieval units reduce the burden on the retriever while the reader can exploit extended contexts to synthesize information from multiple retrieved passages. This hybrid design captures the best of both worlds---efficient retrieval over massive corpora and the reasoning depth of long-context models.

\textbf{Complementary hybrid frameworks} demonstrate that RAG is not a standalone solution but can be integrated with other knowledge infusion techniques to create systems that robustly handle diverse requirements. For instance, RAG can supply verified factual evidence as demonstrations for in-context learning, thereby grounding the model's reasoning in external sources. Conversely, knowledge editing can be applied to the retriever's index to rapidly correct obsolete entries without retraining, while the generator remains frozen. ERAGent \cite{arxiv:2405.06683v1} exemplifies such integration by coupling retrieval with an experiential learner that incrementally updates a user profile and an external knowledge bank, combining the dynamism of RAG with continual, lightweight adaptation. Similarly, memory-augmented models such as MemoRAG \cite{arxiv:2409.05591v2} maintain a global memory of the knowledge base, merging long-term parametric memory with on-demand retrieval. These compound architectures illustrate that the future lies not in choosing one paradigm over another, but in designing modular, semi-parametric systems that fluidly select among parametric recall, non-parametric lookup, and targeted updates according to the specific knowledge requirements and latency constraints of the task.


\section{Retrieval Mechanisms and Knowledge Source Construction}



\subsection{Retrieval Model Architectures and Training Paradigms}


The design of the retrieval model fundamentally shapes the quality and efficiency of a RAG system by determining how documents are represented and matched against queries. The landscape spans sparse lexical methods, dense neural encoders, late-interaction frameworks, and the emerging paradigm of generative retrieval, each with distinct trade-offs in effectiveness, latency, and training complexity.

Classical sparse approaches such as BM25 rely on exact term matching and inverted indices, offering high efficiency and interpretability. Learned sparse representations like SPLADE extend this by using pretrained language models to predict token-level importance, preserving the inverted index while capturing term weighting signals from context, as surveyed in \cite{arxiv:2211.14876v1}. These methods remain competitive for keyword-centric queries but struggle with semantic gaps. In contrast, dense dual-encoder models independently encode queries and documents into fixed-dimensional vectors, enabling fast approximate nearest-neighbor search. Early systems such as DPR \cite{arxiv:2211.14876v1} employ a bi-encoder trained with contrastive objectives that pull relevant pairs close and push irrelevant ones apart. The success of such models depends critically on hard-negative mining strategies that select challenging non-relevant passages, as well as on large-scale pretraining with tasks like masked language modeling or contrastive learning \cite{arxiv:2211.14876v1}. More expressive late-interaction architectures, typified by ColBERTv2, compute fine-grained token-level similarities but defer the costly interaction to a late scoring stage, achieving higher relevance precision while retaining much of the efficiency of dual-encoders \cite{arxiv:2211.14876v1}.

To combine complementary strengths, hybrid systems blend sparse and dense signals via score interpolation, cascade reranking, or learnable fusion, significantly improving robustness for long-tail and ambiguous queries. Such ensembles are particularly effective when individual models exhibit distinct failure modes.

A radical departure is generative retrieval, where the retriever directly generates document identifiers or passage content rather than conducting a vector search. Architectures like DSI and Self-Retrieval treat the index as a set of model parameters and decode document IDs autoregressively from the query, effectively blurring the boundary between retrieval and generation \cite{arxiv:2404.14851v3}. Unified embedder-generator models, such as LLM-Embedder, extend this idea by using a single LLM to handle both retrieval and generation, promising seamless knowledge access but introducing challenges in scalability and latency.

Training these models to align with downstream generation goals is an active frontier. Standard dense retrievers are trained with InfoNCE-type contrastive losses, often augmented with knowledge distillation from cross-encoder teachers \cite{arxiv:2211.14876v1}. To bridge the retriever-generator gap, joint training paradigms that propagate generation-phase reward signals back to the retriever have proven essential. For instance, REPLUG uses the probability improvement of a frozen black-box LM as a learning signal to update the retriever \cite{arxiv:2301.12652v4}, while Atlas jointly pretrains the retriever and LM end-to-end, enabling the entire system to be optimized for answer likelihood \cite{arxiv:2208.03299v3}. Self-RAG takes this further by training a single LM to adaptively decide when to retrieve and how to critique its own output via reflection tokens, merging retrieval decision and generation into a unified process \cite{arxiv:2310.11511v1}. Additionally, instruction tuning of retrieval models, as in INTERS and FollowIR, aligns query and document representations with task-specific user intentions, improving relevance in complex scenarios \cite{arxiv:2401.06532v2,arxiv:2403.15246v1}.

In summary, retrieval model design is moving from static, off-the-shelf components toward deeply integrated, task-aware modules trained jointly with the generator. Future work will likely explore more efficient generative retrieval at scale, better handling of multi-modal corpora, and training objectives that explicitly account for the generator's downstream faithfulness and noise robustness.


\subsection{Indexing and Scalable Retrieval Infrastructures}


Building on the retrieval model advances that demand billion-scale, low-latency search, the scalability of RAG hinges on indexing structures, compression, and system-level orchestration that go far beyond brute-force vector search. Modern dense retrievers project queries and documents into high-dimensional embedding spaces, making exact nearest-neighbor lookups computationally prohibitive. Approximate nearest-neighbor (ANN) indices thus serve as the backbone, trading negligible drops in accuracy for orders-of-magnitude speed-ups. Inverted file systems with product quantization (IVF-PQ) partition the vector space via k-means, store residual quantized codes, and probe only a subset of clusters. Graph-based approaches like hierarchical navigable small-world (HNSW) graphs construct multi-layered proximity structures that allow efficient greedy traversal, achieving higher recall at the cost of increased memory and construction time. The choice of index type defines the recall--latency Pareto frontier and must be aligned with embedding dimension, query volume, and hardware constraints.

To alleviate storage and bandwidth bottlenecks, aggressive compression is essential. Product quantization (PQ) decomposes embeddings into sub-vectors, quantizes each subspace, and approximates inner products with precomputed lookup tables, drastically shrinking memory footprint while retaining semantic fidelity. Binary hashing converts vectors to compact binary codes, enabling ultra-fast Hamming distance comparisons often accelerated by SIMD instructions on CPUs. With the rise of retrieval-augmented small language models deployed on edge devices, quantization extends to both embeddings and model weights at INT4 or INT8 levels, though careful calibration is needed to avoid retrieval quality degradation.

System-level co-design further amplifies throughput by deeply integrating retrieval with generation. PipeRAG \cite{arxiv:2403.05676v1} introduces pipeline parallelism that overlaps retrieval and autoregressive decoding, adaptively scheduling retrieval intervals based on generation state and a learned performance model, delivering up to a 2.6$\times$ end-to-end speed-up without quality loss. Speculative retrieval frameworks such as RaLMSpec \cite{arxiv:2401.14021v1} anticipate future retrieval needs, batch speculative queries, and verify results asynchronously, effectively hiding retrieval latency behind the generative process. Query-vector and result caches exploit temporal locality, while modular toolkits like FlashRAG \cite{arxiv:2405.13576v1} provide pre-built components that streamline the deployment of optimized retrieval pipelines.

Real-world corpora are dynamic, requiring incremental index updates for streaming documents. Online k-means clustering adapts centroids on the fly, and partial re-indexing rebuilds only affected segments, preserving freshness without full index invalidation. Generative retrieval models, including CorpusBrain \cite{arxiv:2402.16767v1} and PAG \cite{arxiv:2404.14600v1}, present a radical alternative by encoding corpus knowledge directly into model parameters, eliminating the external index entirely. While promising for moderate collections, scaling such end-to-end document identifier generation to web-scale corpora remains an open challenge. Together, these infrastructure innovations furnish the high-recall, low-latency search that makes dense and hybrid RAG systems feasible at scale, yet the ultimate effectiveness of these indices is only as good as the queries that probe them---motivating the query processing and augmentation strategies detailed next.


\subsection{Query Processing and Augmentation Strategies}


In retrieval-augmented generation (RAG), the raw user input---often an ambiguous natural language utterance, a multi-turn dialogue fragment, or a complex information need---rarely maps directly to an effective retrieval query. The gap between conversational or task-level expressions and the indexing vocabulary of the external knowledge store necessitates dedicated query processing and augmentation strategies. This subsection systematically reviews techniques that transform inputs into retrieval-ready representations, spanning classical pseudo-relevance feedback, generative expansion, decomposition, and adaptive rewriting, all aimed at maximizing the relevance and coverage of retrieved evidence.

\textbf{Generative query expansion} has emerged as a powerful paradigm, leveraging large language models (LLMs) to enrich a query with related terms, concepts, or even synthetic passages. Unlike traditional pseudo-relevance feedback that relies on top-ranked documents, generative approaches produce diverse contextual expansions without external reliance. For instance, the Generation-Augmented Retrieval \cite{arxiv:2009.08553v4} framework demonstrates that heuristically generated relevant contexts substantially enrich query semantics, enabling BM25 to match or surpass dense retrievers on open-domain QA, and further fusing multiple expansions yields consistent gains. In tool retrieval, Re-Invoke \cite{arxiv:2408.01875v2} generates synthetic queries covering the query space of each tool document and employs intent-based multi-view similarity ranking to pinpoint relevant tools. Domain-specific strategies also exploit structured knowledge: CBR-RAG \cite{arxiv:2404.04302v1} augments legal queries with similar past cases using case-based reasoning, while ERAGent \cite{arxiv:2405.06683v1} incorporates an Enhanced Question Rewriter and Knowledge Filter to iteratively refine the query for complex personalization scenarios.

For multi-hop or compositional information needs, \textbf{sub-query decomposition} translates the original question into a set of simpler retrieval tasks. Corrective Retrieval Augmented Generation \cite{arxiv:2401.15884v2} incorporates a decompose-then-recompose algorithm that selectively focuses on key information from retrieved documents and can trigger web search when confidence is low. Iter-RetGen \cite{arxiv:2305.15294v2} synergizes generation and retrieval in an iterative loop: the partially generated output serves as an informative context to refine the next retrieval, progressively accumulating evidence. More explicit planning-based decomposition is exemplified by Learning to Plan for Retrieval-Augmented Large Language Models from Knowledge Graphs \cite{arxiv:2406.14282v1}, which fine-tunes LLMs on planning data derived from KGs to break complex questions into manageable retrieval steps.

Conversational settings demand \textbf{context-aware reformulation} to handle coreference, ellipsis, and topic shifts. Standalone rewritten queries that capture the full intent (``generate-then-retrieve'') prevent context drift and improve retrieval precision. ChatRetriever \cite{arxiv:2404.13556v1} adapts LLMs for conversational dense retrieval by learning robust session representations through contrastive learning and masked instruction tuning, achieving state-of-the-art performance without explicit rewriting. Meanwhile, instruction-aware and adaptive rewriting uses model confidence or self-reflection to decide when to rewrite. For example, CRAG's lightweight retrieval evaluator assesses document quality and can trigger corrective actions, including query refinement. Such adaptive mechanisms trade off extra computation against retrieval precision, a crucial balance for real-world deployment.

In summary, query processing in RAG has evolved from simple keyword matching to a nuanced orchestration of generative expansion, decomposition, and context-aware rewriting. Current trends point toward deeper integration of LLMs as both the rewriter and the retriever, creating a feedback loop where generation informs retrieval and vice versa. Future work must address the efficiency of these augmentation steps, develop rigorous metrics for query quality, and extend robust processing to multimodal and low-resource domains, ensuring that the query bridge between user intent and knowledge stores becomes a transparent and highly optimized component of RAG architectures.


\subsection{Knowledge Source Design, Preprocessing, and Maintenance}


The construction of a high-quality external knowledge repository is foundational to retrieval-augmented generation, directly enabling the query-side strategies discussed in the previous subsection. The retrievability, relevance, and freshness of the evidence exposed to the language model all originate from how the repository is built and maintained. This lifecycle begins with document parsing and chunking, where raw formats such as PDF and HTML are converted into plain text and segmented into indexable units. The choice of chunk granularity---sentence, passage, or a sliding window---directly affects context density and the risk of introducing noise or truncating semantically coherent information. Industry experience reported by Barnett et al. \cite{arxiv:2401.05856v1} underscores that improper chunking and failure to preserve structural boundaries rank among the most common sources of RAG system degradation. To mitigate such issues, hierarchical indexing strategies are increasingly adopted, employing a two-level structure in which document-level summaries guide retrieval toward relevant sections, and fine-grained passage indices provide precise evidence. More sophisticated approaches augment this with graph-based indices that capture cross-chunk relations, allowing retrieval to span the original document's discourse structure even after fragmentation.

Beyond structural considerations, the semantic alignment between the index and anticipated user queries is critical for retrieval effectiveness---forming the same challenge that query processing techniques seek to bridge from the user's side. This has motivated generative data preparation techniques that precompute metadata to close the gap between raw document content and natural language queries. For example, Promptagator \cite{arxiv:2209.11755v1} demonstrates the power of using a large language model to generate synthetic query--document pairs from a handful of examples, effectively creating task-specific training data for retrievers. Extending this philosophy, recent frameworks construct per-passage meta-knowledge artifacts---such as concise summaries, question--answer pairs, or entity descriptions---that are stored alongside the original text. These enriched representations enable the retriever to match on conceptual intent rather than lexical overlap alone, substantially improving recall in long-tail and multi-faceted information needs.

Modern knowledge sources are increasingly heterogeneous, encompassing images, tables, code snippets, and structured knowledge graphs. Dedicated multimodal retrieval-augmented generators, exemplified by MuRAG \cite{arxiv:2210.02928v2}, establish a unified non-parametric memory that indexes both text and visual modalities, enabling a single generator to reason across evidence types. The integration of knowledge graphs has given rise to GraphRAG systems \cite{arxiv:2408.08921v2}, where entities and relationships are indexed using graph structures, allowing retrieval to exploit topological semantics that complement textual similarity. Such unified architectures demand careful representation alignment so that multimodal and structured data can be seamlessly injected into the language model's context without loss of fidelity.

Sustaining the knowledge repository over time introduces challenges of versioning, temporal awareness, and conflict resolution. In dynamic domains, the retriever must distinguish between outdated and authoritative versions of facts. TempRALM \cite{arxiv:2401.13222v2} explicitly incorporates temporal signals into the retrieval process, combining semantic relevance with recency to answer time-sensitive queries, and demonstrates up to a 74\% performance improvement over temporally agnostic baselines. Equally important is the handling of factual contradictions and redundant content across multiple sources. When multiple passages present divergent claims, the generator risks being misled or overloading the context with contradictory evidence. Methods that cluster semantically equivalent documents and apply evidence synthesis---sometimes referred to as corrective retrieval augmentation---prune noise and highlight consensus, as evaluated in benchmarks that test noise robustness and negative rejection \cite{arxiv:2309.01431v2}. Future directions point toward self-maintaining corpora that incrementally index streaming data, enforce temporal validity windows, and employ automated fact-verification pipelines to resolve conflicts, thereby closing the loop between knowledge ingestion and trustworthy generation.


\section{Integration Strategies for Augmented Generation}



\subsection{Input-Level Context Integration and Prompt Engineering}


Input-level context integration remains the most direct paradigm for fusing retrieved knowledge with generative models, where evidence is prepended to the prompt as additional textual context. This approach, embodied by the original RAG formula \cite{arxiv:2005.11401v4} and later adopted in black-box settings \cite{arxiv:2301.12652v4}, treats the LLM as a static consumer of an augmented input. Despite its simplicity, the performance of this scheme is highly sensitive to the quantity, ordering, and formatting of the injected passages, as well as the inherent attention biases of the underlying language model. Consequently, careful engineering of the prompt structure is essential to balance the benefits of retrieved evidence against the risks of increased context length, noise, and positional dilution.

A primary concern is the management of the context window. Simple concatenation of all retrieved passages can easily exceed the model's input limits, forcing truncation heuristics that may discard crucial information. To mitigate this, systems often adopt parallel encoding followed by sparse attention selection, as in Sparse RAG, or increase retrieval granularity to reduce the number of units while preserving contextual integrity \cite{arxiv:2406.15319v3}. However, even within a fixed window, the \emph{position} of evidence profoundly influences generation quality. The well-known ``lost in the middle'' phenomenon---where information placed in the central portion of a long prompt is under-utilized---motivates explicit document ordering strategies. While reinforcement-based reordering methods (e.g., R$^{4}$) learn to place critical documents at positions of high attention, a more immediately practical lesson emerges from studies on retrieval noise: adding irrelevant documents at specific positions can unexpectedly boost performance \cite{arxiv:2401.14887v3}. These findings underscore that input-level augmentation must be treated as a structured learning problem, not a trivial concatenation.

Beyond raw passage placement, the design of the prompt template plays a decisive role. Delimiting retrieved content with special tokens (e.g., \texttt{<read> ... </read>}) and providing clear task instructions help the model distinguish between external knowledge and its own parametric memory, reducing interference. Moreover, incorporating few-shot in-context demonstrations that model how to synthesise evidence from retrieved documents can substantially improve factual grounding. Self-reflective frameworks further extend prompt engineering by injecting special reflection tokens that allow the model to dynamically decide whether to trust retrieved passages or to abstain from using them \cite{arxiv:2310.11511v1}. Although these approaches move toward adaptive integration, their core mechanism relies on a carefully crafted input template that encodes both retrieved content and meta-instructions.

To address the burden of excessively long contexts, compression at the input level has become an active research focus. Techniques such as those in CRAG \cite{arxiv:2401.15884v2} cluster and decompose retrieved documents, distilling them into concise summaries or key sentences before concatenation, drastically reducing token count while preserving essential evidential value. Similarly, knowledge filtering via query blending \cite{arxiv:2402.11129v1} can eliminate extraneous passages prior to integration, thereby improving the signal-to-noise ratio. These compression methods effectively trade off some fidelity for computational efficiency, making them particularly appealing in latency-sensitive deployments.

Ultimately, the effectiveness of input-level augmentation hinges on a holistic orchestration of retrieval, ordering, and prompt formatting. Future work should explore dynamic, model-aware ordering that leverages attention patterns to place evidence optimally, and investigate how prompt compression can be made more content-preserving without sacrificing conciseness. The interplay between parametric knowledge and formatted evidence remains an open challenge, but the continued improvement of instruction-tuned retriever--generator interfaces promises to make input-level context integration ever more robust and efficient.


\subsection{Deep Architectural Fusion Mechanisms}


While input-level integration treats retrieved knowledge as raw text prepended to the prompt, the text-as-interface paradigm places fundamental limits on information density and breaks end-to-end differentiability, as discussed in the previous section. A deeper class of methods addresses these limitations by embedding external evidence directly into the internal computational graph of the generative model. These architectural fusion mechanisms enable richer, more parameter-efficient, and often more robust integration by manipulating latent representations, cross-attention patterns, and intermediary activations, moving knowledge injection from the input layer into the model's core processing.

A foundational architectural approach is cross-attention-based conditioning, classically exemplified by the Fusion-in-Decoder (FiD) architecture. In FiD, each retrieved passage is independently encoded, and the decoder jointly attends to all passages through its cross-attention layers, allowing the model to distil relevant evidence at a granular, token-specific level. The inherent computational overhead of this design has motivated innovations such as FiD-Light, which constrains information flow from the encoder to the decoder to improve the latency-effectiveness Pareto frontier without sacrificing performance \cite{arxiv:2209.14290v1}. Extending this cross-attention paradigm to decoder-only architectures, which lack a dedicated cross-attention module, remains an active challenge that has spurred the development of adapter-based and latent-space fusion methods.

Moving beyond conditioning on raw token sequences, a family of modular approaches inserts lightweight, trainable adapter layers into a frozen language model to process and integrate retrieved knowledge vectors. The key insight is to decouple knowledge integration from the original language model parameters, preserving strong generative priors while enabling cost-effective fine-tuning. Parallel to this, virtual token strategies introduce a small set of learnable, continuous embeddings (virtual tokens) that are prepended to the prompt. Fine-tuned on downstream tasks, these soft tokens act as highly compressed instructions that steer the model's self-attention to systematically focus on relevant segments of the retrieved context, effectively re-weighting attention without altering the backbone model.

At the extreme frontier of architectural compression lies the direct fusion of dense retrieval embeddings into the language model's representation space, treating retrieved embeddings not as a source for text but as a distinct modality. The xRAG framework pioneers this by employing a lightweight modality bridge---typically a small projection layer---to map document embeddings directly into the transformer's input token space, eliminating the need for explicit text decoding \cite{arxiv:2405.13792v1}. This achieves extreme compression ratios and reduces overall FLOPs by orders of magnitude, while remarkably matching or exceeding the performance of uncompressed models on knowledge-intensive tasks. Similarly, the DRAG paradigm injects compressed entity or passage embeddings into intermediate layers, creating a non-textual channel of evidence grounding that side-steps the token-length bottleneck entirely.

A critical advantage unifying these deep fusion strategies is the potential for end-to-end differentiability. When the retrieval module's output is a continuous vector rather than discrete text, the entire pipeline---from retrieval encoder to fusion module to generator---can theoretically be optimized via gradient descent. The Stochastic RAG framework formalizes this by recasting the discrete retrieval step as a stochastic sampling-without-replacement process, using a straight-through Gumbel-top-k approximation to allow gradients to propagate through the retrieval step \cite{arxiv:2405.02816v1}. This moves beyond separate, decoupled optimization toward a holistic training paradigm where the retriever learns to produce embeddings that the generator's fusion mechanism is specifically optimized to consume. The central challenge for these architectures remains balancing tight integration with the modular, plug-and-play nature that makes RAG systems practical, agile, and maintainable in dynamic knowledge environments. While the methods described so far embed knowledge statically before generation begins, the retrieval process need not be a one-shot preprocessing step; the following section explores adaptive and iterative paradigms that turn retrieval into an active partner throughout the generation process.


\subsection{Adaptive and Iterative Retrieval-Generation Interactions}


Within adaptive and iterative retrieval-generation paradigms, retrieval is transformed from a one-shot preprocessing step into an active partner that collaborates with the generative model throughout the entire output process. Rather than treating retrieval as a static, pre-determined action, these strategies dynamically decide when and what to retrieve, gather evidence across multiple reasoning steps, and incorporate self-reflective critique loops to produce more accurate and grounded long-form outputs.

A foundational component is \textbf{dynamic retrieval triggering} guided by generation confidence or knowledge boundary estimation. Methods such as Corrective Retrieval Augmented Generation \cite{arxiv:2401.15884v2} introduce a lightweight retrieval evaluator that assesses the relevance of initially retrieved documents and returns a confidence score; low confidence triggers alternative actions such as web search or query reformulation, while selective decomposition `recomposes' the retrieved content to filter out noise. Similarly, ERAGent \cite{arxiv:2405.06683v1} incorporates a Retrieval Trigger module that curtails unnecessary external retrieval when internal parametric knowledge suffices, thus balancing efficiency and grounding. The Time-Aware Adaptive Retrieval approach \cite{arxiv:2402.16457v1} further demonstrates that LLMs can be guided to make reliable retrieval decisions without calibration, using temporal cues to decide when the internally stored knowledge becomes outdated. These confidence-gated mechanisms effectively prevent over-retrieval and the introduction of irrelevant context while preserving the ability to tap into external knowledge when the model's own knowledge is insufficient.

Beyond single-step retrieval, \textbf{multi-hop and recursive evidence accumulation} capture the interplay between retrieval and generation in an iterative fashion. Iter-RetGen \cite{arxiv:2305.15294v2} proposes an iterative synergy where the model's generated output serves as a query to retrieve additional relevant knowledge, which then refines the next generation round, processing all retrieved information holistically without structural constraints. IRCoT \cite{arxiv:2212.10509v2} interleaves retrieval with chain-of-thought reasoning, using each generated reasoning step to formulate a targeted retrieval operation; the retrieved evidence in turn informs subsequent reasoning, creating a tight feedback loop that significantly improves multi-step question answering. Such architectures demonstrate that splitting complex information needs into iterative, sub-goal-oriented retrieval steps not only increases factual grounding but also reduces the risk of cascading errors from early incorrect retrievals.

Another critical direction involves \textbf{self-reflective and metacognitive loops} that allow the system to criticise and refine its own outputs. RA-ISF \cite{arxiv:2403.06840v1} decomposes tasks into three submodules that iteratively generate, evaluate, and refine answers through self-feedback, enabling the model to identify gaps and re-retrieve information without external supervision. MemoRAG \cite{arxiv:2409.05591v2} adopts a dual-system design where a lightweight long-range LLM forms a global memory of the corpus and produces a draft answer; this draft then serves as a cluing signal to guide the retrieval process for a second, more expressive LLM that generates the final answer. This memory-inspired architecture effectively performs a form of metacognitive planning, bridging the gap between ambiguous information needs and structured retrieval.

Finally, the synergy between retrieval and explicit reasoning or task planning is increasingly important. By deriving planning data from knowledge graphs and fine-tuning LLMs to construct step-wise retrieval plans \cite{arxiv:2406.14282v1}, complex queries can be decomposed into sub-tasks, each triggering targeted retrieval, and the collected evidence is then consolidated for final generation. Active retrieval strategies in code generation \cite{arxiv:2402.12317v1} similarly employ iterative query refinement and knowledge soup updates, demonstrating the general applicability of the paradigm beyond pure text.

In summary, adaptive and iterative retrieval-generation interactions greatly extend the capability of RAG systems to handle tasks requiring deep reasoning, evolving knowledge needs, and self-correction. Future work will likely pursue tighter integration of long-term memory, more efficient metacognitive loops that dynamically trade off retrieval cost against accuracy, and unified frameworks that seamlessly blend planning, retrieval, and generation into a single, coherent cognitive process.


\subsection{Robust Evidence Selection and Handling Noisy, Conflicting, or Redundant Content}


Building on the preceding discussion, the dynamic retrieval strategies that interleave generation with adaptive and iterative lookups amplify the risk of incorporating irrelevant, contradictory, or redundant evidence. Therefore, faithful integration of retrieved content demands robust evidence selection and noise handling to safeguard the reliability of retrieval-augmented generation (RAG). This subsection surveys mechanisms that span inference-time filtering, training-time robustness enhancement, multi-document evidence synthesis, and uncertainty-aware decoding, each designed to ensure that the generator is guided only by useful, consistent signals from the external corpus.

At inference time, a first line of defense involves filtering or re-ranking the retrieved candidate set to eliminate low-quality passages. Lightweight natural language inference (NLI) scorers and fact-checking modules---such as those integrated into RETA-LLM \cite{arxiv:2306.05212v1}---can post-process retrieval results, discarding documents that are entailed by the query or that fail to provide supporting evidence. Re-ranking pipelines that prioritize documents based on predicted utility or answerability further reduce noise exposure. However, the relationship between relevance and generation improvement is not monotonic; recent work has demonstrated that irrelevant documents can, in some scenarios, unexpectedly boost downstream accuracy \cite{arxiv:2401.14887v3}, underscoring the need for context-dependent filtering that accounts for the complex interaction between retrieved signals and the generator's parametric knowledge. The RGB benchmark \cite{arxiv:2309.01431v2} reveals that many large language models (LLMs) exhibit poor negative rejection---failing to ignore irrelevant content---and struggle with contradictory evidence, making explicit post-hoc filtering a critical safeguard.

Beyond inference-time heuristics, training-time interventions can instill inherent robustness. Methods that expose the model to diverse retrieval noise during fine-tuning enable it to learn to discriminate between useful and misleading context solely from final-answer supervision, as proposed in the RAAT paradigm. While detailed implementations of RAAT lie outside the papers provided here, the necessity of such training is clearly motivated by the RGB benchmark's findings on information integration and counterfactual robustness \cite{arxiv:2309.01431v2}. End-to-end optimization approaches, such as Stochastic RAG \cite{arxiv:2405.02816v1}, align the retriever and generator through differentiable sampling, encouraging the entire pipeline to focus on passages that truly contribute to generation quality and implicitly discouraging reliance on spurious or noisy evidence. These joint training strategies help bridge the gap between retrieval relevance and the generator's eventual needs.

When multiple retrieved passages contain conflicting facts, the generator must synthesize a coherent response or appropriately indicate uncertainty. Multi-hop question-answering benchmarks like MultiHop-RAG \cite{arxiv:2401.15391v1} exemplify the challenge of integrating complementary and sometimes contradictory evidence across documents. Iterative retrieval-generation frameworks, such as Iter-RetGen \cite{arxiv:2305.15294v2}, allow the model to gather additional information based on partial outputs, progressively reconciling conflicts and refining the answer. Nevertheless, explicit conflict detection and resolution remain largely unsolved; current systems often rely on the LLM's implicit reasoning to weigh evidence, which can be unreliable without dedicated contradiction-aware training data or architecture.

Uncertainty-aware decoding provides a final safeguarding layer by monitoring faithfulness in real time and adjusting the generation process accordingly. Metrics like those provided by RAGAS \cite{arxiv:2309.15217v1} can be internalized to measure the semantic alignment between the generated text and the retrieved evidence, enabling beam-search strategies that penalize ungrounded continuations. Lightweight synchrony monitors that analyze decoding dynamics---such as token likelihood and context influence---can flag and correct hallucinations before they are surfaced to the end user. This aligns with the operational insight from practical deployments that validation must be continuous during inference \cite{arxiv:2401.05856v1}.

Taken together, robust evidence selection and noise handling demand a multi-layered approach that combines filtering, training, synthesis, and decoding. The field is moving toward tightly integrated frameworks where these components are jointly optimized, yet significant gaps remain: formalizing the conditions under which noise becomes beneficial, designing annotator-free conflict resolution strategies, and creating faithful monitors that work across diverse domains. Future research must also develop training protocols that explicitly teach models when to trust external evidence versus their parametric knowledge, ultimately yielding RAG systems that remain reliable even when retrieval quality is imperfect.


\section{Training and Optimization of Retrieval-Augmented Systems}



\subsection{End-to-End Joint Training and Differentiable Retrieval}


End-to-end joint training of the retriever and generator constitutes a paradigm shift from modular, independently optimized components to a tightly coupled system where the generation loss can directly shape retrieval behavior. This subsection examines approaches that bridge the retrieval--generation boundary, covering early fine-tuning schemes that marginalize over retrieved passages, reinforcement learning strategies that treat the language model as a black-box reward, and recent differentiable approximations that relax discrete sampling to enable true end-to-end gradient flow.

Foundational work on the RAG model \cite{arxiv:2005.11401v4} introduced joint fine-tuning by treating the retrieved documents as latent variables. The training objective marginalizes over the top-\(k\) passages, \(P(y|\mathbf{x}) = \sum_{\mathbf{z}} P(y|\mathbf{x},\mathbf{z}) P(\mathbf{z}|\mathbf{x})\), and backpropagates through the generator while updating the retriever's query encoder. Although the document index remains fixed and the marginalization is approximated over a small set, this method established the viability of aligning retrieval with generation. Atlas \cite{arxiv:2208.03299v3} scaled this idea to few-shot pre-training by adopting an asynchronous re-indexing mechanism: after several training steps, the entire document index is re-embedded using the updated query encoder, preserving consistency between the retriever and the knowledge store. Such marginalization-based approaches work well when the retriever can be accessed and updated, but they are inherently limited by the top-\(k\) truncation and the heavy cost of periodic re-indexing.

When the generator is a black-box model---common when leveraging proprietary large language models---gradient-based joint training is infeasible. Reinforcement learning provides an alternative supervision channel for the retriever. REPLUG \cite{arxiv:2301.12652v4} treats the language model as an opaque reward function: it computes the probability improvement of the gold sequence when a particular document is prepended and uses this signal to train the retriever via policy gradients. By framing retrieval as an action that maximizes the downstream LM's likelihood, REPLUG aligns the retriever with generation without requiring internal model access. This black-box paradigm sacrifices the fine-grained gradient coupling of end-to-end methods but offers practical utility for augmenting closed-source models and has inspired subsequent approaches that mine hard examples from generator outputs to further refine retrieval.

A more recent thrust eliminates the discrete non-differentiability of top-\(k\) retrieval altogether through continuous relaxations. Stochastic RAG \cite{arxiv:2405.02816v1} recasts retrieval as a stochastic sampling-without-replacement process and employs straight-through Gumbel-top-\(k\) as a differentiable approximation. This enables gradients to flow from the generator's loss back through the sampling operation to the retriever, permitting true end-to-end optimization without marginalization or independence assumptions. Experiments across open-domain QA, fact verification, and dialogue tasks show that such fully differentiable joint training yields significant performance gains, validating the idea that end-to-end gradient descent can harmonize retriever and generator more effectively than two-stage or reward-based approaches.

Complementary to the above, iterative and expectation-maximization-style schemes cyclically update the retriever using generator-derived pseudo-labels and then retrain the generator with improved retrieval, gradually tightening their synergy. While not fully differentiable, these alternating protocols often strike a pragmatic balance between training stability and effectiveness.

Collectively, these methods illustrate a clear trajectory toward deeper integration. Marginalization-based fine-tuning offers simplicity, RL-based supervision provides black-box flexibility, and differentiable retrieval unlocks the full potential of gradient-based learning across the entire RAG pipeline. Remaining challenges include scaling differentiable retrieval to billion-scale indices without prohibitive compute, mitigating the variance introduced by stochastic sampling, and combining joint training with parameter-efficient adaptation to democratize its benefits. Future advances will likely amalgamate these techniques, enabling more coherent, self-improving retrieval-augmented systems that continuously align the retriever with the evolving needs of the generator.


\subsection{Fine-Tuning, Instruction Tuning, and Domain Adaptation}


Building on the joint training paradigms described previously, practical deployment of RAG systems demands adaptation strategies that teach pre-trained LLMs and retrievers to effectively integrate external, often noisy, non-parametric evidence. This subsection examines supervised fine-tuning, dual instruction tuning, and domain-specific customization, with a focus on parameter efficiency, synthetic data augmentation, and mitigating catastrophic forgetting.

Supervised fine-tuning (SFT) forms the bedrock of RAG adaptation. Foundational work by Lewis et al. \cite{arxiv:2005.11401v4} established a recipe where both retriever and generator are fine-tuned on pairs of queries and ground-truth answers, with the generator conditioning on retrieved documents. This trains the LLM to marginalize over disparate passages, learning implicit relevance weighting. The Atlas model \cite{arxiv:2208.03299v3} demonstrated that this joint SFT enables powerful few-shot learning by teaching the model to attend to the most pertinent evidence. However, a persistent challenge is bridging the preference gap between the retriever's notion of relevance, often derived from semantic similarity, and the LLM's utility for generation. To address this, the dual instruction tuning framework RA-DIT performs a two-stage protocol: first updating the LLM with retrieval-augmented instructions to better exploit context, and then fine-tuning the retriever to match the newly aligned LLM's preferences, significantly improving synergy.

As full-parameter fine-tuning of billion-parameter LLMs becomes prohibitively expensive, parameter-efficient fine-tuning (PEFT) has become dominant. Techniques like Low-Rank Adaptation (LoRA) are implicitly or explicitly employed to keep base LLM parameters frozen while injecting lightweight, trainable adapter modules that condition the generation on retrieved knowledge. The AAR approach \cite{arxiv:2305.17331v1} extends this plug-and-play philosophy, training a generic retriever that adapts to frozen, unseen LLMs, which is crucial for black-box and cost-sensitive settings. A critical risk in domain adaptation is catastrophic forgetting of general linguistic and reasoning capabilities. Methods like RE-Adapt tackle this by reverse-engineering an instruction-following adapter from the domain-specialized model, allowing the system to maintain broad competencies while excelling in vertical applications such as medicine, as reported in clinical RAG case studies \cite{arxiv:2402.01733v1}.

The scarcity of domain-specific, retrieval-augmented training data is a significant bottleneck that synthetic data generation aims to resolve. Frameworks like InFO-RAG \cite{arxiv:2402.18150v1} take an unsupervised approach by treating the LLM as an ``Information Refiner,'' training it to produce outputs that are more accurate and concise than the potentially noisy retrieved passages, without requiring paired ground-truth data. Conversely, RADA uses external datasets to retrieve related instances and diversify the fine-tuning corpus. Predetermined retrieval-augmented strategies can also be synthesized; the GenGround framework \cite{arxiv:2406.14891v2} uses a generate-then-ground distillation process to teach smaller models complex multi-step reasoning paradigms.

In summary, these adaptation strategies point toward more integrated, data-efficient, and forgiving approaches to RAG, with future directions including tighter coupling in dual-tower optimization via hindsight posterior-guided training \cite{arxiv:2110.07752v2} and the development of lifelong learning systems that seamlessly blend synthetic data with on-the-fly PEFT to continuously adapt from minimal examples. While these fine-tuning recipes equip generators to leverage external context, the resulting models often still struggle with the nuanced demands of faithful generation, preference alignment, and robustness to imperfect retrieval. The subsequent subsection builds on these adapted architectures and training procedures to examine learning paradigms that couple preference optimization with noise resilience, moving beyond static fine-tuning toward dynamic alignment of generator behavior.


\subsection{Alignment, Preference Optimization, and Noise Robustness Training}


Training retrieval-augmented generators to faithfully integrate external evidence while aligning with human preferences and withstanding retrieval noise is essential for reliable deployment. This subsection examines learning paradigms that couple preference optimization with robustness, moving beyond static fine-tuning toward dynamic alignment of generator behavior with the retrieved context.

Reinforcement learning from human feedback (RLHF) and its AI-feedback variants reward generations that are factually consistent with retrieved sources, penalizing unsupported claims. In the RAG setting, tailoring such signals to the retrieval interface is critical. For example, BIDER \cite{arxiv:2402.12174v1} first distills retrieved documents into key supporting evidence via supervised fine-tuning, then employs reinforcement learning to align the distilled evidence with the LLM's preferences, effectively filtering noise while preserving essential information. This preference-based refinement not only improves answer quality but also reduces input length, demonstrating how alignment objectives can simultaneously address robustness. More broadly, iterative self-feedback frameworks, such as RA-ISF \cite{arxiv:2403.06840v1}, decompose tasks and self-critique using retrieved knowledge, implicitly training the model to recognise and correct inconsistencies through repeated cycles of generation and evaluation.

Contrastive and knowledge-distillation objectives further shape generator behavior. Information-refinement losses teach the LLM to produce answers that are more concise and accurate than the raw retrieved texts, encouraging the model to act as an information refiner. Meanwhile, evaluation-token mechanisms, as in UniMS-RAG \cite{arxiv:2401.13256v1}, embed relevance and consistency scores into the generation process, enabling the model to self-regulate its reliance on evidence. These techniques enable fine-grained control over how external knowledge is used, addressing both faithfulness and personalization demands.

Adversarial training for noise resilience has emerged as a vital component of robust RAG training. Standard fine-tuning often assumes high-quality retrieval, yet real-world retrievers return a mixture of relevant, irrelevant, and contradictory content. To combat this, retrieval-augmented adaptive adversarial training (RAAT) dynamically exposes the generator to varied noise types---irrelevant passages, conflicting documents, and redundancy---while only providing final-answer supervision. This multi-task learning regime forces the model to internalize noise-recognition heuristics, preserving generation quality even under severe retrieval degradation. Notably, systematic analyses have revealed that the presence of irrelevant documents can sometimes improve accuracy in certain configurations, underscoring the need to design training curricula that exploit such counter-intuitive benefits rather than simply filtering noise \cite{arxiv:2401.14887v3}.

Addressing knowledge conflicts remains a frontier challenge. Preference optimization techniques can be extended to teach the generator when to trust parametric knowledge over retrieved context and vice versa. By constructing datasets with deliberate contradictions and applying reinforcement learning or direct preference optimization, models learn to arbitrate between conflicting sources, mitigating both contextual blindness and over-reliance on retrieval. BIDER's approach of refining evidence to match LLM preferences exemplifies how alignment objectives can be repurposed for conflict resolution \cite{arxiv:2402.12174v1}.

Looking forward, the integration of alignment, preference optimization, and robustness training will likely coalesce into unified meta-learning frameworks that adapt to corpus dynamics and user feedback continuously. Combining introspection loops with adversarial noise injection could yield self-improving RAG systems that maintain factual integrity across diverse and evolving knowledge landscapes. The development of lightweight, sample-efficient preference pipelines and the formalisation of robustness-alignment trade-offs will be central to advancing this promising direction.


\subsection{Efficient Inference and Deployment Optimization}


Serving these dependable generators at scale introduces unique computational bottlenecks that transcend those of standalone generation, primarily stemming from the overhead of processing and conditioning on large volumes of retrieved text. Addressing these challenges calls for a holistic co-design of algorithms and systems. A foundational concern is the management of key-value (KV) caches, which store intermediate attention states to avoid recomputation. In RAG pipelines, retrieved documents often share common passages across queries or iterative retrieval rounds, creating opportunities for reuse. RAGCache introduces a knowledge tree structure that organizes the KV states of documents in a hierarchical cache across GPU and host memory. By employing replacement policies tailored to the access patterns of RAG---where documents are retrieved and then attentively read---this approach dramatically reduces the time-to-first-token and overall memory pressure, transforming a na\"{i}ve bottleneck into a reusable, high-speed library of knowledge states.

Beyond caching, speculative decoding has been adapted to exploit the deterministic nature of provided context. In standard speculative decoding, a draft model proposes multiple future tokens that are verified in parallel by the target LLM. REST reformulates this for RAG by using the retrieved documents themselves as the source for drafting, hypothesizing that the generation will align closely with the retrieved text. This shifts the draft generation from a separate small model to a simple string-matching or copy mechanism from the context. Extending this to iterative retrieval scenarios, RaLMSpec introduces speculative retrieval with batched verification, which parallelizes the retrieval for a future generation step while the current decoding is underway, effectively hiding retrieval latency behind computation. These methods guarantee that the output distribution remains unchanged while achieving substantial speedups, a critical property for maintaining generation quality.

Techniques for prompt compression directly attack the latency and memory costs that scale quadratically with input length in transformer architectures. Sparse RAG encodes retrieved documents in parallel and then employs a learned, query-aware sparsification module during the autoregressive decoding phase. This allows the generator to attend only to the most relevant caches from a potentially vast set of retrieved contexts, transforming a dense attention problem into a highly selective retrieval within the context itself. From a representation perspective, extreme compression methods like xRAG \cite{arxiv:2405.13792v1} fundamentally eliminate the need for textual context altogether. By treating dense retrieval embeddings as a distinct modality and training a lightweight modality bridge, xRAG fuses these embeddings directly into the LLM's representation space, achieving a compression ratio that reduces entire documents to single vectors while reducing FLOPs by a factor of 3.53---a paradigm shift that decouples generation cost from corpus size.

System-level orchestration further amplifies these algorithmic gains through pipeline parallelism. PipeRAG \cite{arxiv:2403.05676v1} formalizes the co-design of retrieval and generation as a performance optimization problem, introducing flexible retrieval intervals to maximize the overlap of sequential processes. Its integrated performance model dynamically balances retrieval quality against latency by adapting the scheduling to both the generation state and the underlying hardware, achieving up to a 2.6$\times$ speedup in end-to-end generation. For resource-constrained environments, these design philosophies extend to edge deployment, where compute-in-memory backed RAG systems explore hardware--software co-design to make on-device retrieval-augmented inference viable. The trajectory of the field points toward a unified optimization framework in which KV-cache management, speculative retrieval, extreme context compression, and hardware-aware parallelism are not isolated techniques but composable primitives within an intelligent inference engine, dynamically reconfigured to meet stringent service-level objectives for latency, throughput, and accuracy while faithfully serving the robust generators shaped by the alignment and noise-resilience training described earlier.


\section{Evaluation, Benchmarking, and Trustworthiness Assessment}



\subsection{Evaluation Dimensions, Metrics, and Attribution}


The evaluation of Retrieval-Augmented Generation (RAG) systems demands a multifaceted framework that transcends the traditional evaluation of isolated retrieval or generation models. It necessitates a holistic assessment that captures the synergy between the retriever and generator, with a central focus on attribution---the degree to which generated content can be traced back to and supported by the retrieved evidence \cite{arxiv:2312.10997v5}. This section establishes the foundational metrics for this assessment, moving from component-level analysis to integrated measures of fidelity.

Component-level evaluation provides the initial diagnostic lens. Retrieval quality is typically gauged through classical Information Retrieval (IR) metrics such as Recall, Precision, and Mean Reciprocal Rank (MRR), which assess the ability of the retriever to surface relevant documents from the knowledge corpus \cite{arxiv:2404.10981v1}. For generation, traditional lexical overlap metrics like Exact Match (EM) and F1 score are giving way to more nuanced semantic similarity measures, including BLEURT and BERTScore, which better correlate with human judgment of response quality \cite{arxiv:2405.07437v2}. However, these isolated metrics are insufficient, as a highly relevant retrieved passage does not guarantee a faithful generation, and a superficially fluent answer may be factually unmoored from its sources \cite{arxiv:2312.10997v5}.

This gap is bridged by integrated attribution and groundedness metrics, which form the cornerstone of RAG-specific evaluation. The core concept is "attributable to identified sources" (AIS) evaluation, which jointly scores the fidelity of a generation to its provided evidence. Frameworks like RAGAs (Retrieval Augmented Generation Assessment) have pioneered reference-free, automated metrics for this purpose, decomposing evaluation into dimensions of context relevance, answer faithfulness (groundedness), and answer relevance \cite{arxiv:2309.15217v1}. Such frameworks treat attribution not as a binary variable but as a nuanced spectrum, often leveraging lightweight fine-tuned language models as judges. Automated suites like ARES extend this by using prediction-powered inference with a minimal set of human annotations to fine-tune LLM judges, achieving high correlation with manual evaluation for context relevance and faithfulness \cite{arxiv:2405.07437v2}. Similarly, the eRAG method evaluates each retrieved document individually by measuring the downstream generation performance when conditioned on it, using the task-specific metric as a direct relevance label for the document, thereby creating a strong correlation between retrieval evaluation and overall system performance \cite{arxiv:2404.13781v1}.

To diagnose fundamental system capabilities beyond aggregate scores, targeted metrics have emerged. The RGB benchmark operationalizes these by evaluating critical LLM behaviors essential for RAG success, including noise robustness (degradation under irrelevant contexts), negative rejection (the ability to abstain from answering when evidence is lacking), information integration (synthesizing data from multiple documents), and counterfactual robustness (adherence to evidence over parametric knowledge) \cite{arxiv:2309.01431v2}. This diagnostic approach moves beyond treating a RAG model as a monolith, instead pinpointing specific failure modes such as an over-reliance on parametric knowledge or an inability to resolve conflicts between sources \cite{arxiv:2307.11019v2}. The ultimate frontier is the development of unified evaluation ecosystems that integrate these component, fidelity, and diagnostic metrics, providing a reproducible and comprehensive profile of a RAG system's trustworthiness and moving from simple performance snapshots to an understanding of the conditions under which a system remains reliable.


\subsection{Benchmark Ecosystems, Evaluation Frameworks, and Datasets}


Building on the foundational metrics and diagnostic approaches discussed above, the evaluation of RAG systems has matured from ad-hoc component-wise assessments on isolated question-answering datasets toward comprehensive, reproducible testbeds that probe the intricate interplay between retrieval quality, generation faithfulness, and end-to-end utility. Early RAG studies relied on classical open-domain QA benchmarks such as Natural Questions and TriviaQA \cite{arxiv:2005.11401v4}, yet these proved insufficient to capture the full spectrum of failure modes, including noise robustness, multi-hop reasoning, and long-form attribution. In response, the community has constructed multifaceted benchmark ecosystems, automated evaluation frameworks, and modular toolkits that jointly enable rigorous, comparable, and diagnostically rich analysis.

A landmark step was the \textbf{KILT} benchmark, which unified five knowledge-intensive tasks---fact checking, entity linking, slot filling, open-domain QA, and dialogue---providing a standardized retrieval corpus (Wikipedia) and shared provenance annotations. By evaluating both retrieval and generation under consistent conditions, KILT forced models to handle diverse evidence requirements and spurred development of retrieval-augmented architectures that could be directly compared. Nevertheless, its tasks remain largely answer-centric; more recent benchmarks address specialized domains and broader application scenarios. For code generation, \textbf{CodeRAG-Bench} \cite{arxiv:2406.14497v1} systematically categorizes programming tasks (basic, open-domain, repository-level) and aggregates retrieval sources from competition solutions, documentation, and GitHub repositories, revealing that current retrievers often falter due to limited lexical overlap and that generators cannot reliably exploit long or multiple contexts. Similarly, the trend toward personalized dialogue and unstructured document analysis is driving the creation of domain-specific test collections that stress-test the capacity to synthesize information across heterogeneous sources.

Complementing static datasets, \textbf{automated evaluation frameworks} have emerged to decouple system assessment from costly human annotation. \textbf{ARES} \cite{arxiv:2311.09476v2} generates synthetic training data to fine-tune lightweight language-model judges for context relevance, answer faithfulness, and answer relevance, employing prediction-powered inference with a minimal set of human ratings to achieve high correlation with manual evaluation. \textbf{RAGAS} \cite{arxiv:2309.15217v1} introduced a suite of reference-free metrics that separately quantify retrieval precision, context utilization, and generation faithfulness, enabling rapid development cycles without ground-truth answers. For a more fine-grained retrieval diagnosis, \textbf{eRAG} \cite{arxiv:2404.13781v1} evaluates each retrieved document by feeding it individually to the generator and measuring downstream task performance; this document-level utility signal correlates better with end-to-end RAG performance than conventional relevance labels and opens the door to constructing more informative retrieval training objectives.

The need for fair, reproducible comparison has also spurred the creation of \textbf{modular research toolkits} that standardize the implementation and benchmarking of diverse RAG algorithms. \textbf{FlashRAG} \cite{arxiv:2405.13576v1} offers an open-source library pre-implementing 12 advanced RAG methods and 32 benchmark datasets within a unified, customizable framework, drastically reducing the effort to reproduce, compare, and develop new techniques. Similarly, \textbf{RETA-LLM} \cite{arxiv:2306.05212v1} provides plug-and-play modules for request rewriting, passage extraction, answer generation, and fact verification, enabling systematic exploration of the retrieval-generation interface. These toolkits embody a shift toward ``evaluation-as-a-service'' that prioritizes transparency, extensibility, and diagnostic logging of failure points---exactly the kind of infrastructure envisioned in experience reports that identify seven common failure modes \cite{arxiv:2401.05856v1}.

These integrated evaluation ecosystems and modular toolkits now provide the diagnostic granularity needed to systematically probe a wide spectrum of RAG quality dimensions, from retrieval precision to complex reasoning fidelity and attribution. One dimension that consistently emerges from this infrastructure as a linchpin for trustworthy generation is factual consistency---the precise alignment of generated statements with the retrieved evidence. The next section therefore dives into dedicated detection and mitigation techniques for ensuring this consistency, building directly on the evaluation capabilities outlined here to fortify RAG systems against hallucination and unfaithful generation.


\subsection{Factuality, Faithfulness, and Hallucination Detection and Mitigation}


Ensuring the factual consistency of generated outputs against retrieved evidence stands as a central challenge in RAG systems, motivating a diverse array of automated detection and mitigation strategies. Hallucination detection methods, which identify when generated content diverges from its grounding context, have evolved from simple reference-based verification to more sophisticated self-consistency probes and decoding-time interventions. A foundational class of detectors operates at the level of fine-grained claim verification, decomposing generated responses into atomic assertions for independent validation against the retrieved context. For instance, \cite{arxiv:2309.15217v1} and \cite{arxiv:2310.01558v1} exemplify this approach by checking each claim at a triplet level. These methods often employ Natural Language Inference (NLI) models as automatic evaluators to determine if each atomic claim is entailed by the retrieved passages, a technique foundational to reference-free evaluation suites like RAGAS \cite{arxiv:2309.15217v1}. One critical insight from these approaches is that overall answer accuracy is often less informative than the proportion of perfectly supported sub-claims, highlighting the need for more granular attribution metrics.

Building on this, a parallel family of methods explores self-reflective and zero-resource hallucination detection. These techniques leverage the LLM's internal knowledge and reasoning capabilities to introspect on its own outputs without requiring external references. The Chain-of-Verification framework, for example, uses a multi-step process where an LLM generates verification questions for its initial draft and then independently answers them to check for inconsistency. Similarly, consistency-based methods query the model about the content of its own generated citations to see if it can accurately recall the source material, effectively revealing when information was fabricated. These self-reflective loops are powerful for detecting subtle hallucinations that are not obvious from a single evaluation pass but introduce significant computational overhead due to multiple inference calls.

The stark trade-off between detection accuracy, latency, and cost has driven the development of specialized, cost-effective detection models. While general-purpose LLMs can serve as high-quality evaluators, their deployment at scale is prohibitively expensive. Research has shown that compact, fine-tuned encoder models can match or even exceed the performance of massive LLMs for this specific task. A prime example is Luna, a fine-tuned DeBERTa-large model that achieves higher accuracy than GPT-3.5 for RAG hallucination detection while reducing costs by 97\% and latency by 91\%. This trend underscores a crucial insight for production systems: specialized, lightweight classifiers trained on synthesized hallucination data can provide a Pareto-optimal solution, enabling real-time faithfulness monitoring synchronous with generation.

Beyond detection, mitigation strategies are increasingly integrated directly into the generation process the moment a hallucination risk is identified. Decoding-time techniques represent a proactive paradigm that manipulates output token probabilities to enhance factuality. The Induce-then-Contrast Decoding (ICD) method penalizes logits that would produce induced untruthful continuations, thereby steering the model toward more generation grounded in the context. Similarly, adaptive retrieval frameworks like Rowen use multilingual semantic-aware consistency checks to selectively trigger a secondary retrieval round when internal knowledge gaps suggest an emerging hallucination. These dynamic mitigation strategies are complemented by adaptive RAG architectures such as Self-RAG and CRAG, which use special reflection tokens or lightweight evaluators to decide between parametric generation, on-the-fly retrieval, or rewriting the query \cite{arxiv:2401.15884v2}. In conclusion, the field is converging on an integrated, systems-level approach where lightweight, specialized detectors monitor the generation process in real-time, selectively activating computationally expensive self-reflection or retrieval loops to correct potential factual errors before they reach the end-user, optimizing the balance between trustworthiness and computational feasibility.


\subsection{Robustness, Security, and Privacy Evaluation}


Ensuring trustworthiness under adversarial and non-ideal conditions becomes paramount as retrieval-augmented generation (RAG) moves into high-stakes deployments. This subsection surveys evaluation methodologies that stress-test RAG systems along four tightly intertwined axes: robustness to noisy retrieval, resilience against security attacks, privacy leakage through external databases, and the propagation of bias from retrieved corpora.

Robustness evaluations probe how generators handle irrelevant, misleading, or contradictory evidence. The RGB benchmark \cite{arxiv:2309.01431v2} decomposes this challenge into four diagnostic capabilities: noise robustness, negative rejection (abstaining when evidence is absent), multi-document information integration, and counterfactual robustness. Experiments across six representative LLMs reveal that while models exhibit a degree of noise tolerance, they consistently fail to reject irrelevant context and often integrate conflicting information in harmful ways, leading to factual degradation. A counter-intuitive finding from \cite{arxiv:2401.14887v3} shows that injecting irrelevant documents can, under certain conditions, improve answer accuracy by over 30\%---underscoring that the relationship between retrieval quality and generation fidelity is far from monotonic and that conventional precision-based retrieval metrics may not align with downstream performance. These observations motivate adaptive adversarial training (e.g., RAAT) that exposes the generator to diverse noise types---random irrelevance, subtle contradictions, and redundant passages---in order to harden it against real-world retrieval failures.

The retrieval-generation interface also expands the attack surface beyond traditional text generation. While the reviewed literature focuses heavily on membership inference, the modular nature of RAG invites broader threats such as corpus poisoning and indirect prompt injection. A seminal privacy evaluation, \cite{arxiv:2405.20446v2}, demonstrates a black-box membership inference attack (MIA) that determines whether a specific text passage resides in the external database by querying the RAG system and analyzing the similarity and perplexity of its outputs. The attack achieves high success rates across multiple generative models without gradient access, exposing a fundamental tension between retrieval depth and data confidentiality. The same study proposes initial defense strategies---like instruction-template hardening---that reduce attack success for some datasets, but no universal countermeasure exists yet. These results underscore the urgent need for privacy-preserving retrieval mechanisms, such as on-device dense retrieval or differential privacy at indexing time, as well as integrity-verification layers that detect tampered retrieval results.

Fairness and bias evaluation in RAG remains nascent. Biases embedded in external corpora can be amplified through retrieval and fed directly into the generator's output, potentially skewing decisions in high-stakes domains. Current benchmarks, including RGB, provide indirect signals---for instance, counterfactual robustness testing acts as a proxy for resistance to skewed evidence---but lack dedicated fairness audits that measure demographic or viewpoint biases in generated text as a function of the retrieval source. Moreover, the interplay between parametric knowledge and retrieved context can cause the model to over-rely on a biased snippet when the retriever surfaces it as highly ranked, a phenomenon that warrants controlled experimentation.

In sum, robustness, security, privacy, and fairness are deeply interwoven in RAG systems. The field has made initial progress through diagnostic benchmarks and targeted attack studies, yet substantial gaps remain: characterizing the real-world noise distribution captured by retrieval failures, designing defenses against prompt injection and data extraction attacks, formalizing privacy guarantees for dynamic knowledge bases, and constructing fairness evaluation suites that account for both retrieval and generation biases. Addressing these gaps will be essential for the trustworthy deployment of retrieval-augmented models in open, adversarial, and regulated environments.


\section{Applications and Domain-Specific Adaptations}



\subsection{Knowledge-Intensive Question Answering and Fact Verification}


Knowledge-intensive question answering (QA) and fact verification represent canonical application domains where retrieval-augmented generation (RAG) has proven indispensable, precisely because they demand the grounding of outputs in authoritative external evidence. Early foundational RAG models \cite{arxiv:2005.11401v4} demonstrated that a parametric sequence-to-sequence model, when coupled with a dense retriever over Wikipedia, could outperform both purely parametric models and task-specific extractive architectures on open-domain QA, while also producing more specific and factual language. This semi-parametric paradigm was soon extended to larger scales and few-shot settings: RETRO \cite{arxiv:2112.04426v3} showed that retrieval from a 2-trillion-token database can compensate for a 25$\times$ reduction in parameters, and Atlas \cite{arxiv:2208.03299v3} achieved strong few-shot QA performance by carefully pre-training the retriever and generator jointly. Meanwhile, REPLUG \cite{arxiv:2301.12652v4} established that even black-box LLMs benefit from prepended retrieved documents, and that the LLM itself can supervise the retriever to find more helpful evidence.

A central challenge in these tasks is the potential for retrieval to introduce irrelevant, noisy, or contradictory information that degrades generation quality. Systematic analyses \cite{arxiv:2401.14887v3} have revealed the counter-intuitive phenomenon that irrelevant documents can sometimes improve accuracy, but overall robustness remains fragile. To address this, Corrective Retrieval Augmented Generation (CRAG) \cite{arxiv:2401.15884v2} introduces a lightweight retrieval evaluator that assesses the quality of the initial retrieval and triggers either knowledge refinement via a decompose-then-recompose algorithm or an extended web search when confidence is low. Similarly, BlendFilter \cite{arxiv:2402.11129v1} enhances retrieval relevance by blending internal and external query expansions and then filtering out extraneous information using the LLM's own discriminative capability. On the training side, Retrieval-Augmented Adaptive Adversarial Training (RAAT) \cite{arxiv:2405.20978v1} dynamically exposes the model to diverse noise types and employs multi-task learning to build internal noise-recognition skills, substantially raising F1 and EM scores under adversarial conditions.

Multi-hop reasoning and long-form generation further complicate the picture, as a single-shot retrieve-and-generate strategy often fails to provide all necessary evidence. Iter-RetGen \cite{arxiv:2305.15294v2} processes retrieval and generation in an iterative synergy, using the partially generated output as a query to retrieve additional knowledge, which in turn improves subsequent generation. FLARE \cite{arxiv:2305.06983v2} takes an active, forward-looking approach by predicting upcoming sentences and using low-confidence tokens as triggers for re-retrieval and re-generation, thereby progressively accumulating relevant information throughout the generation of long-form answers. Self-RAG \cite{arxiv:2310.11511v1} elevates adaptivity by training a single LLM to generate special reflection tokens that control on-demand retrieval, critique the relevance of retrieved passages, and assess the factual support of its own output, achieving significant gains in both accuracy and fine-grained citation quality. These dynamic architectures are complemented by Speculative RAG \cite{arxiv:2407.08223v1}, which uses a smaller specialist to draft answers from distinct subsets of documents, while a larger generalist verifies the best draft, thus improving both speed and accuracy.

Fact verification demands not only correct answers but also precise attribution of claims to sources. The LLM-Augmenter system \cite{arxiv:2302.12813v3} iteratively revises responses by generating factuality scores from the retrieved evidence and prompting the model to correct itself. Self-RAG's reflection tokens inherently provide token- or sentence-level provenance, while benchmarks such as RGB \cite{arxiv:2309.01431v2} have diagnosed that current LLMs struggle with negative rejection, information integration, and counterfactual robustness---capabilities essential for reliable fact verification. In response, the field increasingly focuses on metacognitive regulation and selective retrieval. Methods like Rowen \cite{arxiv:2402.10612v1} use multilingual semantic-aware consistency checks to detect potential hallucinations and trigger retrieval only when a knowledge gap is identified, thereby balancing parametric and non-parametric knowledge while preserving efficiency.

Looking forward, the trajectory is toward architectures that seamlessly interweave retrieval, reasoning, and self-assessment in a continuous loop. Challenges remain in handling genuinely contradictory evidence, scaling multi-hop reasoning to open-web corpora, and providing human-verifiable citation chains that go beyond surface-level passage attribution. The synthesis of retrieval-augmented generation with structured knowledge (e.g., knowledge graphs) and the development of more robust, cost-effective fact-verification pipelines will likely define the next generation of trustworthy QA systems.


\subsection{Conversational AI and Dialogue Systems}


Building on the retrieval-augmented generation (RAG) advances in single-shot knowledge-intensive tasks, conversational AI and dialogue systems introduce a distinct set of challenges that demand more dynamic, memory-augmented, and personalized retrieval architectures. Moving beyond isolated QA turns, conversational RAG must handle multi-turn reasoning, user personalization, and mixed-initiative interactions, where underspecified queries, reliance on dialogue history, and the need for dynamic knowledge integration render static retrieval paradigms insufficient. Early work in retrieval-augmented dialogue recognized the value of leveraging similar dialogues as a knowledge source; the skeleton-to-response framework \cite{arxiv:1809.05296v5} showed that revising a retrieved skeletal response can guide more diverse and informative generation than conditioning directly on raw retrieved text. This early insight foreshadowed a core tension: how to mediate between the structured knowledge required for coherent dialogue and the open-ended nature of conversational search.

Modern dialogue RAG systems have evolved along two axes: retrieval-control granularity and memory/personalization sophistication. For retrieval control, a black-box integration paradigm---prepending retrieved passages to the conversational context---has proven effective. RePlug and augmentation-adapted retrievers \cite{arxiv:2305.17331v1} demonstrate that aligning the retriever's output distribution with the LLM's preference, learned from a smaller source LM, can significantly boost factuality in dialogue settings. However, multi-hop queries unfolding over multiple turns necessitate iterative, agent-driven strategies. The Iter-RetGen framework \cite{arxiv:2305.15294v2} loops retrieval and generation so that a partially generated response informs the next retrieval query---a dynamic critical when information needs emerge mid-interaction. IM-RAG \cite{arxiv:2405.13021v1} further advances this by learning an ``inner monologue'' via reinforcement learning, where the LLM Reasoner decides to query the retriever or answer, engaging in a silent multi-round dialogue with the retriever and achieving flexibility across heterogeneous IR systems.

Personalization and long-term memory form the second axis. Conversational agents must ground responses in both external facts and the user's evolving profile and interaction history. MemoRAG \cite{arxiv:2409.05591v2} implements a dual-system architecture: a lightweight, long-range LLM maintains a global memory of dialogue history and user data, generates a draft answer as a clue, and triggers a targeted retrieval round that a more expressive generation LLM then synthesizes. This decoupling enables efficient session-spanning personalization. Similarly, UniMS-RAG \cite{arxiv:2401.13256v1} handles multiple user-specific knowledge bases by having the LLM output special ``acting tokens'' that dynamically select sources, retrieve evidence, and evaluate relevance in a unified sequence-to-sequence loop, iteratively refining the response through self-reflection.

Contemporary systems further enhance dialogue robustness through metacognition and structured knowledge integration. MetaRAG \cite{arxiv:2402.11626v1} introduces a metacognitive regulation pipeline that evaluates the quality of retrieved context and generated response, diagnoses knowledge gaps, and plans corrective re-retrieval cycles---critical for handling conversational ambiguity. GraphRAG-based approaches \cite{arxiv:2408.04948v1} enable precise retrieval of relational facts from knowledge graphs, which is particularly valuable in task-oriented dialogues needing to query structured databases (e.g., customer service ticket relations). Overall, conversational RAG is transitioning from a retrieval-then-generate monolith to a modular cognitive architecture where retrieval, memory, reasoning, and self-assessment cooperate in dynamic loops. These iterative, self-corrective retrieval paradigms, originally honed in dialogue, naturally extend to other domains requiring adaptive, feedback-driven reasoning---a theme we explore next in code-centric tasks. Key remaining challenges include minimizing loop latency for real-time spoken interaction, resolving contradictions between prior turns and newly retrieved evidence, and ethically managing the extensive user profiles that personalized memory architectures accumulate, all while balancing generative fluency with precise, verifiable grounding.


\subsection{Code Generation and Software Engineering}


Retrieval-augmented generation addresses fundamental limitations of parametric language models in code-centric tasks, where APIs, libraries, and domain-specific practices evolve rapidly and internal parametric knowledge quickly becomes outdated. By grounding generation in dynamically retrieved external knowledge---documentation, API specifications, structurally similar code snippets, and historical bug-fix pairs---RAG enhances functional correctness, reduces hallucinations, and supports a wide range of software engineering activities, including code completion, program repair, code summarization, and cross-lingual synthesis \cite{arxiv:2108.11601v2,arxiv:2402.12317v1}.

Early frameworks such as REDCODER \cite{arxiv:2108.11601v2} exemplify the paradigm by extending dense retrieval to construct unimodal (code-only or text-only) and bimodal (code--description) databases. The retriever identifies relevant code or summaries, which are then prepended to the generator's input, leading to substantial gains in both code synthesis and comment generation for Java and Python. The approach demonstrates that properly aligned dense representations can capture functional similarity beyond surface lexical overlap, a critical requirement given that semantically equivalent code fragments often share no tokens. However, such one-shot retrieval may fail when the information need is multifaceted. ARKS \cite{arxiv:2402.12317v1} addresses this through a ``knowledge soup'' that integrates web search, official documentation, execution feedback, and evolved code snippets, combined with an active retrieval strategy that iteratively refines queries using partial outputs and test results. This dynamic multi-source augmentation substantially improves execution accuracy on benchmarks featuring long-tail programming languages and frequently updated libraries, proving the value of execution signals in guiding retrieval.

Automated program repair further benefits from retrieval of relevant bug-fix pairs. Although not explicitly represented in our cited corpus as a standalone work, the REDCODER framework's capacity to fetch code pairs can be repurposed to retrieve patches, while ARKS's feedback loop provides a model for iterative repair where generated patches are validated against test suites and retrieval is re-triggered on failure. The central challenge is bridging the syntactic--semantic gap: conventional text-based retrievers often fail to surface functionally relevant artifacts because code expresses logic through structural patterns rather than natural language. This has motivated the development of code-specific dense encoders and retrieval objectives that incorporate execution semantics, yet even state-of-the-art retrievers struggle on reasoning-intensive code retrieval tasks, as evidenced by the BRIGHT benchmark \cite{arxiv:2407.12883v1}, where leading models achieve merely 18.0 nDCG@10 on programming-related queries. Such results underscore the need for retrieval methods that perform deeper program understanding.

Another critical obstacle is context redundancy and the limited context windows of generators. Retrieved documentation and snippets can be verbose, demanding compression or selective filtering. While extreme compression techniques like xRAG \cite{arxiv:2405.13792v1} have been demonstrated for natural language, their extension to code contexts, where structure must be preserved, remains an open frontier. Latency is equally paramount in developer workflows: speculation and pipeline parallelism, as in PipeRAG \cite{arxiv:2403.05676v1}, can overlap retrieval and generation, while KV-cache offloading via RetrievalAttention \cite{arxiv:2409.10516v2} reduces memory footprint and speeds up token generation, directly applicable to real-time code completion.

Looking forward, tighter integration between retrieval, execution-based verification, and language-model reasoning promises to yield self-improving pipelines that not only recall relevant artifacts but also refine them through iterative compilation and testing. Robust, code-aware retrieval models, coupled with efficient inference strategies, will be instrumental in embedding RAG seamlessly into modern developer environments.


\subsection{High-Stakes and Specialized Domain Applications}


This drive toward retrieval-execution-reasoning loops is equally critical in high-stakes domains such as healthcare, law, finance, and scientific research, where retrieval-augmented generation (RAG) must satisfy exceptionally stringent requirements for factual accuracy, regulatory compliance, and domain-specific terminology. Unlike open-domain scenarios, these applications demand fine-grained provenance, reliable citation of authoritative sources, and robust mitigation of hallucination---often under strict latency and interpretability constraints. Consequently, RAG architectures are being adapted to integrate deeply with specialized knowledge structures, incorporate iterative reasoning, and support expert-in-the-loop validation.

Healthcare and biomedicine have seen the most concentrated RAG customization efforts. The MIRAGE benchmark \cite{arxiv:2402.13178v2} reveals that combining multiple medical corpora with hybrid retrievers can elevate LLM accuracy by up to 18\%, while also exposing losses from the ``lost-in-the-middle'' effect. In clinical deployment, a case study integrating 35 preoperative guidelines via RAG demonstrated that GPT-4 augmented with retrieval achieved 91.4\% accuracy, surpassing junior doctors and demonstrating non-inferiority to human experts while delivering answers orders of magnitude faster \cite{arxiv:2402.01733v1}. To address the need for iterative evidence gathering, systems like IM-RAG use reinforcement-learning-guided multi-round retrieval, where the LLM learns inner monologues to decide when to query further or finalize an answer, a paradigm crucial for differential diagnosis or complex treatment guideline synthesis \cite{arxiv:2405.13021v1}. Failures in biomedical RAG---including missing relevant passages and handling contradictory evidence---are systematically documented in case studies that identify seven recurrent engineering pitfalls, emphasizing the need for continuous validation during operation \cite{arxiv:2401.05856v1}.

Legal and financial domains impose analogous demands but with even stricter citation of statutes, regulations, and financial reports. While few dedicated RAG studies exist for these fields, the core enabling techniques---hybrid retrieval blending sparse and dense signals \cite{arxiv:2404.07220v1}, query rewriting conditioned on task instructions \cite{arxiv:2405.06683v1}, and temporally aware retrieval that respects document versioning \cite{arxiv:2401.13222v2}---are directly transferable. Multi-Head RAG \cite{arxiv:2406.05085v1} is particularly promising for legal queries that require reconciling multiple facets of a case, as it leverages multi-head attention activations to fetch semantically diverse documents. Moreover, the clustering and summarization of redundant retrieved content (e.g., via CRAG-style compression) reduces noise and ensures that only authoritative, non-contradictory context reaches the generator, a critical safeguard for compliance-driven outputs.

Scientific research and education benefit from RAG's ability to handle long, dense corpora and produce instructional content with verifiable grounding. LongRAG \cite{arxiv:2406.15319v3} processes entire documents as 4K-token units, drastically reducing the number of retrieval candidates and achieving competitive QA performance on scientific benchmarks without task-specific training. Memory-augmented paradigms like MemoRAG \cite{arxiv:2409.05591v2} form global representations of a knowledge base, enabling the system to answer ambiguous, exploratory queries typical in literature review or textbook generation. For structured knowledge integration, GraphRAG \cite{arxiv:2408.08921v2} indexes and retrieves over knowledge graphs, capturing entity relationships essential for drug-discovery pipelines or physics concept mapping. Extreme context compression methods such as xRAG \cite{arxiv:2405.13792v1} allow efficient injection of document embeddings directly into the LLM's latent space, reducing computation while preserving retrieval semantics---an advantage for resource-constrained educational platforms.

Looking forward, the path to production-grade high-stakes RAG lies in tighter coupling of retrieval with domain-specific reasoning loops, automated fact-checking modules that operate at the claim level, and proactive compliance monitoring. As these systems mature, they will increasingly embody a synergy between non-parametric knowledge access and rigorous expert workflows, enabling trustworthy AI assistance in the most critical human endeavors.


\section{Open Challenges and Future Research Directions}



\subsection{Toward a Principled Theoretical Understanding of RAG}


While the empirical success of retrieval-augmented generation has been amply demonstrated, the field still lacks a rigorous theoretical framework capable of predicting \emph{when} and \emph{why} augmentation helps or harms output quality. Moving beyond trial-and-error heuristics toward a predictive science of RAG is a critical open challenge. Recent work has begun to formalise the benefit--detriment duality inherent in providing external context. Studies such as \cite{arxiv:2402.13492v3} construct fine-grained, fact-centric analyses that reveal retrieval is not uniformly beneficial; its impact depends on the joint frequency of entities and relations in a query, with larger LMs handling popular facts well but struggling with rare entity-relation pairs where a retriever excels. Complementary findings in \cite{arxiv:2401.14887v3} show that even irrelevant documents can unexpectedly improve performance by over 30\%, underlining the need to mathematically decouple the beneficial and detrimental effects of retrieved passages on generation probability.

A principled understanding must also clarify the interaction between parametric knowledge stored in model weights and non-parametric knowledge accessed via retrieval. Research probing knowledge boundaries, such as \cite{arxiv:2307.11019v2}, demonstrates that LLMs frequently display overconfidence and poorly perceive their own limitations, but that retrieval augmentation can sharpen this boundary awareness. This aligns with work like \cite{arxiv:2402.11457v1}, which proposes that retrieval should be triggered based on the model's own certainty, thereby framing the decision as an adaptive trade-off between internal and external resources. Similarly, \cite{arxiv:2212.10511v4} shows that retrieval-augmented LMs significantly outperform pure parametric models on long-tail knowledge, while parametric knowledge remains competitive for high-popularity facts, establishing a quantitative foundation for when to rely on which memory system.

Translating these insights into design guidelines requires formal models that predict the generator's behaviour when supplied with retrieved context. For example, \cite{arxiv:2310.01558v1} provides a systematic analysis of failure cases and proposes training-time data mixing that teaches the model to discern useful from misleading evidence, a step toward principled robustness. Extending such work through distribution-divergence measures between the generation distributions with and without retrieval, or by quantifying the mutual information between retrieved passages and output tokens, could yield objective functions that optimally balance parametric and non-parametric contributions. Future theoretical efforts should aim to derive guidelines for retrieval granularity, integration depth, and corpus composition that are grounded in the statistical properties of the knowledge store and the model's own epistemic uncertainty. Achieving this will transform RAG from an art of empirical tricks into a predictable, verifiable component of reliable language generation systems.


\subsection{Advancing Multimodal, Cross-Modal, and Structured Knowledge Integration}


Building on the need for a principled, predictive science of RAG outlined earlier, the evolution of retrieval-augmented generation beyond plain text to encompass multimodal and structured knowledge represents a critical frontier for building truly comprehensive AI systems. The vast majority of real-world knowledge is not siloed in unstructured text but is distributed across images, audio, video, tables, and intricate knowledge graphs. A system's ability to answer complex queries---such as ``What was the public reaction captured in viral images after the event described in this report?''---demands seamless reasoning over heterogeneous evidence sources, exactly the kind of challenge where theoretical understanding of when and why retrieval helps becomes most urgent. This subsection examines the key challenges and emerging paradigms for advancing multimodal, cross-modal, and structured knowledge integration in RAG, focusing on cross-modal alignment, robustness to noise, and the design of unified generative architectures.

\textbf{Pioneering Multimodal and Cross-Modal Retrieval}  
Early systems established a bi-modal bridge by treating non-textual data as a complementary knowledge source. Works like MuRAG \cite{arxiv:2210.02928v2} and RA-CM3 demonstrated that accessing an external non-parametric memory containing both images and text significantly boosts performance on knowledge-intensive tasks, achieving up to 10-20\% absolute accuracy improvements. These models leverage pre-trained dual-encoders like CLIP for retrieval and jointly train a generator on a mixture of modalities, thereby enabling novel capabilities such as faithful image generation grounded in external visual context.

However, a core challenge in such pipelines is the \emph{multi-granularity noisy correspondence} (MNC) problem. Retrieval from the wild often yields results that are semantically related but functionally irrelevant or misleading, a problem amplified when aligning noisy image-text pairs. Recent work tackles this through robust reranking strategies. For instance, the RagLLaVA/RagVL framework \cite{arxiv:2407.21439v2} instruction-tunes a multimodal large language model (MLLM) to serve as a knowledge-enhanced reranker, precisely filtering the top-k retrieved images before generation. Furthermore, noise-injected training at both the data and token levels during generator fine-tuning enhances systemic robustness, ensuring that the final output is faithful to the most reliable evidence.

\textbf{Integrating Structured and Graph-Based Knowledge}  
Moving beyond free-form modalities, the integration of structured knowledge from databases and knowledge graphs (KGs) presents a distinct set of opportunities and challenges. Unlike the dense vector spaces of text and images, KGs encode complex relational structures and entity-centric facts that are poorly captured by sequential context. The emerging paradigm of GraphRAG \cite{arxiv:2408.08921v2} formalizes this workflow, employing graph-based indexing to preserve relational dependencies and guide retrieval for more precise, multi-hop reasoning.

A critical architectural insight is the complementarity of vector-based semantic search and graph-based structural traversal. Systems like HybridRAG \cite{arxiv:2408.04948v1} demonstrate that fusing context from both a vector database and a knowledge graph outperforms either approach individually. In a financial information extraction task, HybridRAG improved retrieval accuracy and answer generation quality by jointly leveraging the semantic fluency of vector search and the factual precision of graph querying. This synergy is further evidenced in domain-specific applications; in customer service, amalgamating RAG with a KG \cite{arxiv:2404.17723v2} preserved the critical intra-issue structure and inter-issue relations often lost by text-chunking, reducing median resolution time by 28.6\%.

\textbf{Toward Unified Architectures and Faithful Grounding}  
The ultimate challenge is the design of a single, unified architecture that can fluidly ingest and reason over any combination of retrieved evidence. Efforts like UniMS-RAG \cite{arxiv:2401.13256v1} propose a unified multi-source framework where a sequence-to-sequence paradigm uses special tokens to dynamically select and evaluate knowledge from different sources. The system generates ``acting tokens'' to interact with various knowledge stores and ``evaluation tokens'' to gauge relevance, coupled with a self-refinement mechanism that iteratively improves the response based on consistency and relevance scores.

A parallel line of research seeks extreme compression to bridge the modality gap. The xRAG approach \cite{arxiv:2405.13792v1} reinterprets document embeddings from a dense retriever as a distinct ``retrieval modality,'' fusing them directly into the LLM's representation space via a learned modality bridge. This eliminates the need to expand retrieved documents into long text prompts, achieving a dramatic compression rate and a 3.53$\times$ reduction in FLOPs while outperforming uncompressed models. This philosophy of treating retrieved content from any source as a composable feature vector points toward a future where the generator's context window is a rich, lossless tapestry of multimodal and structured embeddings.

\textbf{Synthesis and Future Directions}  
The path forward demands a three-pronged focus. First, research must move beyond simple concatenation toward deeper, latent-space fusion mechanisms that can balance the influence of conflicting modalities without overwhelming the generator. Second, developing unified data augmentation and adversarial training strategies that simulate realistic multimodal retrieval noise is crucial for hardening systems against cross-modal hallucination. Finally, the evaluation ecosystem must evolve to include benchmarks that explicitly test for compositional reasoning over mixed modalities, going beyond unimodal metrics to assess the attribution fidelity of generated answers to the specific parts of a heterogeneous evidence set. By surmounting these challenges, the next generation of RAG systems will serve as universal information hubs, capable of leveraging the full spectrum of the world's knowledge. Realizing this vision sets the stage for the next critical requirement: the management of temporal evolution across these diverse, dynamic knowledge sources, a parallel challenge addressed in the following subsection.


\subsection{Continual Knowledge Updating and Lifelong Learning in RAG}


While retrieval-augmented generation decouples knowledge storage from model parameters, enabling direct access to fresh information, this architecture introduces a new critical challenge: maintaining the synchronization, coherence, and utilization efficiency of the non-parametric memory over time without catastrophic forgetting or prohibitive computational cost. We identify three interconnected axes of this challenge. First, the \textbf{temporal integrity and incremental evolution of the knowledge base} itself must be managed. Second, \textbf{retrieval models must engage in continual learning} to adapt their representations to shifting data distributions. Third, the overall system must support \textbf{lifelong learning and self-improvement} by leveraging its own generated outputs and interaction history as new, curated memory.

A foundational requirement is the construction of dynamic, temporally aware indexing infrastructures. Real-world deployments in news, finance, and software development demand knowledge bases that ingest streaming data and serve time-sensitive queries. This necessitates moving beyond static vector indices to support incremental updates, versioned snapshots, and conflict resolution across document versions. The practical implications of this are highlighted in code generation and software engineering contexts, where frequently updated libraries render static knowledge obsolete; systems like \cite{arxiv:2402.12317v1} actively combine web search and documentation to construct a dynamic "knowledge soup." A failure to maintain temporal awareness can lead to the retrieval of contradictory or outdated evidence, directly undermining the factual grounding that RAG promises, as observed in case studies of RAG system failure points \cite{arxiv:2401.05856v1}.

The second axis involves the adaptation of the retrieval model itself, which confronts the classic stability--plasticity dilemma. Dense retrievers, often fine-tuned on static corpora, can suffer severe performance degradation when the target document distribution shifts. Continual pre-training or sequential fine-tuning of the retriever on new data risks catastrophic forgetting of previously learned indexing structures. While rehearsal-based methods and elastic weight consolidation offer partial solutions, the open challenge lies in developing retrieval-aware continual learning strategies that can efficiently update the embedding space without full re-indexing. This is particularly critical for domain-specific retrievers, such as those in biomedicine \cite{arxiv:2404.18443v1}, which must track evolving medical literature. Unsupervised adaptation techniques that leverage the target corpus for retriever selection, as in \cite{arxiv:2402.04853v1}, represent a pragmatic step, but fall short of true, sustained online learning.

The third, most forward-looking axis, concerns self-improving and lifelong RAG systems. Here, the system's own successful reasoning traces, problem-solving steps, and final outputs are stored, managed, and reused as a form of synthetic, experiential memory. Frameworks like \cite{arxiv:2305.02437v3} exemplify the core mechanism, where an iterative loop uses a retrieval-augmented generator to create an unbounded memory pool and a memory selector to curate the most valuable output for future use, effectively enabling the model to bootstrap its own performance. This paradigm is extended in agent-centric architectures that incorporate an "Experiential Learner" to incrementally model user profiles and expand knowledge \cite{arxiv:2405.06683v1}. The shift toward memory-inspired RAG, as in \cite{arxiv:2409.05591v2}, uses a lightweight long-context model to form a global memory, cluing retrieval tools for later tasks, thereby enabling systems to handle ambiguous information needs by drawing on a compressed representation of past experiences.

Synthesizing these directions, the future of continual RAG lies in developing unified, metacognitive architectures. Such systems would dynamically orchestrate time-sensitive retrieval from versioned corpora, continuously fine-tune their retrieval components against distribution shift, and curate a structured memory of past interactions. This involves moving from passive, single-shot retrieval to an active, self-reflective process where the model not only consumes information but also strategically writes back to its own long-term memory, progressively refining its knowledge and reasoning capabilities over its operational lifetime.


\subsection{Privacy-Preserving, Secure, and Trustworthy RAG Systems}


As RAG systems mature with temporally-aware indexing and self-improving memory, their deployment in high-stakes and regulated domains such as healthcare, finance, and law brings to the fore a critical tension between the benefits of external knowledge access and the risks to privacy, security, and trustworthiness. This subsection examines these emergent vulnerabilities and the nascent body of research dedicated to mitigation, establishing a foundational trustworthiness layer for building responsible RAG systems.

A primary privacy threat is the potential for the external retrieval corpus to contain sensitive information, making it a target for extraction. Anderson et al. \cite{arxiv:2405.20446v2} directly demonstrated this vulnerability through black-box Membership Inference Attacks (MIAs), whereby an adversary can determine with high confidence whether a specific text passage exists in the RAG's database by simply analyzing the model's output distribution. This finding fundamentally challenges the assumption that retrieval databases are opaque to end-users. A core research direction for privacy-preserving RAG involves moving computation to the client side, with architectures like on-device dense retrieval toolkits that process sensitive corpora locally. Complementary strategies include applying differential privacy during index construction or generating synthetic, privacy-compliant data that retains the utility of the original corpus for retrieval while masking individual identities. These approaches must be calibrated to balance the privacy-utility trade-off, ensuring that sanitization does not render the retrieved context ineffective for downstream generation.

Beyond privacy leakage, the retrieval-generation interface introduces a novel attack surface for adversarial manipulation. Cuconasu et al. \cite{arxiv:2401.14887v3} empirically revealed the counter-intuitive finding that strategically injecting irrelevant documents can significantly alter RAG output, a phenomenon that extends to malicious corpus poisoning. An attacker could inject crafted passages to trigger backdoor behaviors or manipulate the model's generated opinions on sensitive topics. Securing RAG against such threats requires robust, multi-layered defenses. This includes developing integrity-verification layers that cryptographically sign trusted documents, adversarial training regimes that harden the generator against misleading or poisoned contexts, and inference-time filtering methods using Natural Language Inference (NLI) scorers to detect and discard non-entailing or spurious passages. The challenge is to maintain the flexibility of open-corpus retrieval while establishing a verifiable chain of trust in the provenance and integrity of the retrieved evidence.

Ensuring trustworthiness ultimately demands mechanisms for verifiable faithfulness and provenance. While RAG is designed to ground generation in evidence, the system can still hallucinate or misattribute information. This has motivated the development of automated evaluation frameworks like RAGAS \cite{arxiv:2309.15217v1} and diagnostic benchmarks such as RGB \cite{arxiv:2309.01431v2}, which explicitly test critical capabilities like noise robustness and negative rejection. Moving beyond evaluation, operational trustworthiness is predicated on real-time faithfulness monitors. Techniques like SynCheck analyze decoding dynamics to detect unsupported claims on the fly, serving as a guardrail against subtle contextual hallucinations. Furthermore, the creation of tamper-evident retrieval logs is essential for high-stakes domains, providing a transparent, auditable trail from the generated output back to the exact source documents, thereby enabling post-hoc verification and fostering accountability. The future of this area lies in the co-design of privacy, security, and faithfulness mechanisms into a unified trustworthiness layer, evolving RAG from a performance-enhancing paradigm into a reliable infrastructure for knowledge-critical applications.


\subsection{Autonomous Agents, Tool Use, and Compositional Reasoning with RAG}


The transition of large language models from static knowledge responders to autonomous agents capable of executing long-horizon tasks demands a fundamental reimagining of retrieval-augmented generation (RAG). Rather than a one-shot prefix of context, retrieval becomes a proactive, tool-like action that an agent dynamically invokes to plan, reason, and interact with the world. This paradigm shift requires seamlessly integrating retrieval into agentic loops that support multi-step planning, adaptive tool use, and metacognitive self-reflection.

At the core of agent-driven retrieval are dynamic strategies that let the model decide \emph{when} and \emph{what} to retrieve based on the evolving generation state. Self-RAG \cite{arxiv:2310.11511v1} equips a language model with trainable reflection tokens that signal on-demand retrieval, relevance filtering, and output critique, enabling a single model to adaptively switch between parametric and non-parametric knowledge. FLARE \cite{arxiv:2305.06983v2} takes a forward-looking approach, using low-confidence token predictions to anticipate future content and trigger explicit retrieval, thus interleaving generation and information seeking in long-form outputs. Corrective Retrieval Augmented Generation (CRAG) \cite{arxiv:2401.15884v2} introduces a lightweight retrieval evaluator that decides if retrieved documents suffice, triggers a web-search fallback when confidence is low, and applies decompose-then-recompose algorithms to eliminate noise---thereby casting retrieval as a composable, quality-aware tool. Iter-RetGen \cite{arxiv:2305.15294v2} demonstrates that iteratively using partially generated output as a query for the next retrieval round yields stronger relevance and multi-hop accuracy, while Adaptive-RAG \cite{arxiv:2403.14403v2} unifies these strategies by learning a complexity-driven router that selects among no-retrieval, single-step, or iterative retrieval paths, effectively planning the retrieval flow.

Beyond pure retrieval, autonomous agents must orchestrate a heterogeneous set of tools---databases, APIs, code interpreters---and reason over their outputs. The modular RAG framework \cite{arxiv:2407.21059v1} provides a blueprint for such architectures by decomposing RAG into independent, recombinable modules with routing, scheduling, and fusion operators, enabling design patterns like conditional and looping retrieval that are essential for task-oriented assistants. ERAGent \cite{arxiv:2405.06683v1} exemplifies this principle by coupling a retrieval trigger with a knowledge filter and an experiential learner that incrementally models user profiles, thereby personalizing multi-turn interactions and minimizing redundant retrieval. These systems illustrate a trajectory where RAG is not merely a prepended context but a decision-making component that can be called, skipped, or chained according to an explicit plan.

Memory and metacognition further elevate agentic RAG for long-horizon tasks. Self-reflective loops, as in RA-ISF \cite{arxiv:2403.06840v1}, iteratively decompose a problem and refine answers through self-feedback, closely mimicking human metacognitive checking. Explicit memory modules, such as MEMORYLLM \cite{arxiv:2402.04624v1} and MemLLM \cite{arxiv:2404.11672v1}, allow agents to maintain a persistent, updatable knowledge store that can be read and written during an episode, bridging the gap between short-term retrieval context and lifelong learning. The ability to store compressed memories of past successful reasoning chains, as hinted by approaches that learn from feedback \cite{arxiv:2112.09737v2,arxiv:2206.06520v1}, further enables agents to self-improve over repeated interactions without full retraining.

Finally, compositional reasoning over heterogeneous knowledge sources is critical when agents must combine facts from structured knowledge graphs, unstructured text, and multimodal documents. Benchmarks like MultiHop-RAG \cite{arxiv:2401.15391v1} reveal that current systems struggle to chain evidence from multiple documents. Hybrid approaches such as HybridRAG \cite{arxiv:2408.04948v1} that fuse vector-based retrieval with graph traversal point toward a future where the retrieval tool itself can query different indices through a unified plan. Ultimately, the path forward lies in end-to-end learnable agents that treat retrieval, tool invocation, and memory operations as differentiable actions within a compositional reasoning framework, enabling robust, verifiable, and general-purpose intelligence.


\subsection{Efficiency, Sustainability, and Democratization of RAG}


As agentic loops multiply the computational demands of retrieval and demand rigorous, context-diverse testing, the widespread adoption of Retrieval-Augmented Generation (RAG) is currently constrained by steep computational costs, a considerable carbon footprint, and a lack of standardized evaluation frameworks that genuinely reflect diverse deployment contexts and linguistic backgrounds. Alleviating these practical barriers is pivotal for democratizing the technology, enabling its transition from well-provisioned cloud environments to edge devices and under-represented language communities. Progress requires coordinated advances across three axes: lightweight model design and inference acceleration, hardware-algorithm co-optimization for on-device deployment, and the construction of open, reproducible, and culturally inclusive benchmark ecosystems.

A central strategy for reducing resource demands involves extreme compression of the retrieval-to-generation pathway. The xRAG framework reinterprets dense retrieval embeddings as features from a distinct modality and fuses them directly into the language model's hidden space via a lightweight modality bridge, completely avoiding textual expansion of retrieved passages and slashing floating-point operations by a factor of 3.53 while maintaining or even improving task accuracy \cite{arxiv:2405.13792v1}. On the retrieval front, the PLAID engine harnesses centroid-based pruning within a highly optimized late-interaction pipeline, reducing search latency by up to 45$\times$ on CPUs without sacrificing retrieval quality \cite{arxiv:2205.09707v1}. These algorithmic gains are complemented by knowledge distillation and careful model-scaling analyses: smaller distilled models can approach the in-domain performance of their larger counterparts, yet parameter count and early query-document interaction remain crucial for zero-shot generalization, indicating that efficiency must be traded off against robustness \cite{arxiv:2206.02873v5}. Striking the right balance is essential for deployment in resource-constrained settings.

A complementary, more radical paradigm is generative retrieval, which encodes the entire corpus within the model parameters and forgoes explicit vector indices. The CorpusBrain model illustrates how strong pre-training tasks can yield a generative retriever that dramatically simplifies the search pipeline and reduces memory footprint \cite{arxiv:2208.07652v1}. Nevertheless, scaling such approaches to millions of passages remains unresolved, with current evidence showing that straightforward parameter scaling is insufficient and that synthetic document representations are vital \cite{arxiv:2305.11841v1}. Further acceleration can be achieved through system-level innovations such as key-value cache reuse across retrieval rounds, speculative decoding adapted for RAG, and pipeline parallelism that overlaps retrieval with generation; all require tight algorithm-system co-design to realize their full potential.

Looking toward edge deployments, future RAG systems must exploit computing-in-memory, browser-based vector search, and hybrid cloud-client architectures to deliver low-latency, private, and energy-efficient augmentation on consumer devices. Such designs will be critical in healthcare, personal assistants, and low-connectivity environments where off-device data transfer is infeasible.

Democratization equally depends on open, reproducible evaluation infrastructures that encompass diverse linguistic and cultural contexts. The FlashRAG toolkit exemplifies a modular, open-source framework that implements a dozen advanced RAG methods over 32 benchmarks, enabling fair, head-to-head comparisons and accelerating research iteration \cite{arxiv:2405.13576v1}. Reference-free evaluation suites such as RAGAS automate the assessment of retrieval relevance and answer faithfulness, drastically reducing the need for costly human annotations \cite{arxiv:2309.15217v1}. Moreover, the generation of synthetic test collections via large language models offers scalable, budget-friendly pathways for creating evaluation resources that can be adapted to new domains and languages \cite{arxiv:2405.07767v1}. Despite these advances, existing benchmarks predominantly focus on English and a few knowledge-intensive tasks. Future efforts must prioritize multi-turn conversational RAG, low-resource languages, and culturally grounded corpora, extending evaluation to include robustness, fairness, and instruction-following capabilities to ensure that the benefits of retrieval augmentation are globally equitable.

In summary, efficiency and sustainability in RAG demand a holistic co-design of algorithms, systems, and evaluation methodology. Lightweight compression, generative retrieval, and hardware-accelerated inference will drastically reduce the environmental footprint, while inclusive, transparent benchmarking will guarantee that progress is broadly beneficial. Together, these directions will transform RAG from a computationally expensive luxury into a sustainable, universally accessible technology, unlocking the agentic vision at global scale.


\section{Conclusion}


As this survey has comprehensively delineated, Retrieval-Augmented Generation (RAG) has rapidly matured from a nascent architectural proposition into a foundational paradigm for building reliable, knowledgeable, and verifiable large language model (LLM)-based systems. The foundational insight, crystallized by Lewis et al. \cite{arxiv:2005.11401v4}, that decoupling parametric memory from a non-parametric, easily updatable external knowledge store can directly address the intrinsic limitations of static LLMs, has proven to be a profound conceptual pivot. This decoupling directly mitigates the critical challenges of factual hallucination, temporal knowledge cutoffs, and the prohibitive cost of domain adaptation via retraining \cite{arxiv:2112.04426v3,arxiv:2208.03299v3}. Our systematic journey through the RAG landscape reveals a field that has successfully consolidated its core techniques while simultaneously expanding into a rich ecosystem of specialized strategies and multifaceted applications.

A primary insight from our analysis is the dialectical progression from monolithic, one-shot augmentation to highly dynamic, adaptive, and self-reflective systems. The initial paradigm of single-round retrieval and input-level concatenation, while powerful in its simplicity and black-box applicability \cite{arxiv:2301.12652v4}, has been decisively transcended. Contemporary research extensively explores iterative retrieval-generation loops and adaptive strategies where the model actively decides when and what to retrieve, significantly enhancing factual grounding for complex, long-form generation tasks \cite{arxiv:2305.06983v2,arxiv:2305.15294v2}. This evolution is epitomized by frameworks like Self-RAG \cite{arxiv:2310.11511v1}, which embed metacognitive critique and on-demand retrieval directly into the generative process via a trainable single model, and by systems that incorporate corrective mechanisms to query diverse external sources when initial retrieval quality is found lacking \cite{arxiv:2401.15884v2}. These advances represent a fundamental shift from treating retrieval as a preprocessing step to integrating it as a core, executable capability of the agent.

Concurrently, the field has developed a sophisticated understanding of the fragility introduced by the retrieval interface and has engineered robust countermeasures. We have detailed the emergence of an entire sub-domain dedicated to managing noisy, irrelevant, conflicting, and redundant evidence, which can lead to a previously underestimated new class of external hallucinations \cite{arxiv:2310.01558v1,arxiv:2405.20978v1}. The development of inference-time filtering via Natural Language Inference, credibility-aware re-ranking, and training-time robustness enhancements like adversarial data augmentation collectively form a critical resilience layer that safeguards the generator from being misled by imperfect retrieval components---a scenario inevitable in real-world deployment \cite{arxiv:2310.01558v1,arxiv:2405.20978v1}. The refinement of evaluation frameworks has kept pace, moving beyond coarse-grained end-task metrics toward fine-grained, component-level diagnostics that specifically measure noise robustness, negative rejection, and information integration capabilities, thereby enabling more principled system development and diagnosis \cite{arxiv:2309.01431v2,arxiv:2309.15217v1}.

Nevertheless, significant theoretical and practical chasms remain, which define the immediate research frontier. A prominent open question concerns the optimal coordination point between retrieval augmentation and the expanding long-context capabilities of next-generation LLMs \cite{arxiv:2407.16833v1}. While RAG offers unmatched computational efficiency and scalable access to vast, evolving corpora, long-context models eliminate the retrieval latency and information fragmentation inherent in chunked indexing. The future is likely hybrid, where systems intelligently route queries based on complexity or utilize retrieved documents as a precision-attended subset within a larger context window, a direction already presaged by self-routing mechanisms. Equally pressing is the need for a formal, principled theory that predicts when retrieval will be beneficial, detrimental, or superfluous, moving the field from an empirically driven practice to a predictive science \cite{arxiv:2402.13492v3,arxiv:2402.11457v1}. Such a theory must account for the intricate interplay between parametric knowledge, retrieved context, and the reasoning capabilities of the generator, ultimately guiding the principled construction of future RAG architectures.

In conclusion, retrieval-augmented generation has irrevocably altered the trajectory of LLM development. It has successfully operationalized the semi-parametric ideal, providing a scalable and tractable pathway to grounding language model outputs in verifiable external data. Our survey charts the paradigm's progression from the canonical "retrieve-then-read" pipeline to a rich design space of multi-modal, continual, agentic, and deeply adaptive systems \cite{arxiv:2210.02928v2,arxiv:2408.08921v2}. The central role of RAG in building trustworthy, attributable, and perpetually up-to-date AI systems is now unequivocally established. As the field pivots toward autonomous agents, lifelong learning, and compositional reasoning, the foundational principles of dynamic knowledge integration, critical evidence assessment, and robust noise management that have been painstakingly developed within the RAG community will constitute the essential cognitive architecture for the next generation of artificial intelligence.

\begin{thebibliography}{101}
\setlength{\itemsep}{2pt}

\bibitem{arxiv:2005.11401v4}
Patrick Lewis, Ethan Perez, Aleksandra Piktus, Fabio Petroni, Vladimir Karpukhin, Naman Goyal \emph{et al.}.
\newblock \emph{Retrieval-Augmented Generation for Knowledge-Intensive NLP Tasks}.
\newblock arXiv:2005.11401v4, 2020.
\newblock \url{https://arxiv.org/abs/2005.11401v4}

\bibitem{arxiv:2401.01313v3}
S. M Towhidul Islam Tonmoy, S M Mehedi Zaman, Vinija Jain, Anku Rani, Vipula Rawte, Aman Chadha \emph{et al.}.
\newblock \emph{A Comprehensive Survey of Hallucination Mitigation Techniques in Large Language Models}.
\newblock arXiv:2401.01313v3, 2024.
\newblock \url{https://arxiv.org/abs/2401.01313v3}

\bibitem{arxiv:2212.10511v4}
Alex Mallen, Akari Asai, Victor Zhong, Rajarshi Das, Daniel Khashabi, Hannaneh Hajishirzi.
\newblock \emph{When Not to Trust Language Models Investigating Effectiveness of Parametric and Non-Parametric Memories}.
\newblock arXiv:2212.10511v4, 2022.
\newblock \url{https://arxiv.org/abs/2212.10511v4}

\bibitem{arxiv:2310.07521v3}
Cunxiang Wang, Xiaoze Liu, Yuanhao Yue, Xiangru Tang, Tianhang Zhang, Cheng Jiayang \emph{et al.}.
\newblock \emph{Survey on Factuality in Large Language Models Knowledge, Retrieval and Domain-Specificity}.
\newblock arXiv:2310.07521v3, 2023.
\newblock \url{https://arxiv.org/abs/2310.07521v3}

\bibitem{arxiv:2312.05934v3}
Oded Ovadia, Menachem Brief, Moshik Mishaeli, Oren Elisha.
\newblock \emph{Fine-Tuning or Retrieval Comparing Knowledge Injection in LLMs}.
\newblock arXiv:2312.05934v3, 2023.
\newblock \url{https://arxiv.org/abs/2312.05934v3}

\bibitem{arxiv:2402.01364v2}
Tongtong Wu, Linhao Luo, Yuan-Fang Li, Shirui Pan, Thuy-Trang Vu, Gholamreza Haffari.
\newblock \emph{Continual Learning for Large Language Models A Survey}.
\newblock arXiv:2402.01364v2, 2024.
\newblock \url{https://arxiv.org/abs/2402.01364v2}

\bibitem{arxiv:2402.11457v1}
Shiyu Ni, Keping Bi, Jiafeng Guo, Xueqi Cheng.
\newblock \emph{When Do LLMs Need Retrieval Augmentation Mitigating LLMs' Overconfidence Helps Retrieval Augmentation}.
\newblock arXiv:2402.11457v1, 2024.
\newblock \url{https://arxiv.org/abs/2402.11457v1}

\bibitem{arxiv:2402.04624v1}
Yu Wang, Xiusi Chen, Jingbo Shang, Julian McAuley.
\newblock \emph{MEMORYLLM Towards Self-Updatable Large Language Models}.
\newblock arXiv:2402.04624v1, 2024.
\newblock \url{https://arxiv.org/abs/2402.04624v1}

\bibitem{arxiv:2202.01110v2}
Huayang Li, Yixuan Su, Deng Cai, Yan Wang, Lemao Liu.
\newblock \emph{A Survey on Retrieval-Augmented Text Generation}.
\newblock arXiv:2202.01110v2, 2022.
\newblock \url{https://arxiv.org/abs/2202.01110v2}

\bibitem{arxiv:2112.04426v3}
Sebastian Borgeaud, Arthur Mensch, Jordan Hoffmann, Trevor Cai, Eliza Rutherford, Katie Millican \emph{et al.}.
\newblock \emph{Improving language models by retrieving from trillions of tokens}.
\newblock arXiv:2112.04426v3, 2021.
\newblock \url{https://arxiv.org/abs/2112.04426v3}

\bibitem{arxiv:2211.14876v1}
Wayne Xin Zhao, Jing Liu, Ruiyang Ren, Ji-Rong Wen.
\newblock \emph{Dense Text Retrieval based on Pretrained Language Models A Survey}.
\newblock arXiv:2211.14876v1, 2022.
\newblock \url{https://arxiv.org/abs/2211.14876v1}

\bibitem{arxiv:2301.12652v4}
Weijia Shi, Sewon Min, Michihiro Yasunaga, Minjoon Seo, Rich James, Mike Lewis \emph{et al.}.
\newblock \emph{REPLUG Retrieval-Augmented Black-Box Language Models}.
\newblock arXiv:2301.12652v4, 2023.
\newblock \url{https://arxiv.org/abs/2301.12652v4}

\bibitem{arxiv:2310.11511v1}
Akari Asai, Zeqiu Wu, Yizhong Wang, Avirup Sil, Hannaneh Hajishirzi.
\newblock \emph{Self-RAG Learning to Retrieve, Generate, and Critique through Self-Reflection}.
\newblock arXiv:2310.11511v1, 2023.
\newblock \url{https://arxiv.org/abs/2310.11511v1}

\bibitem{arxiv:2405.02816v1}
Hamed Zamani, Michael Bendersky.
\newblock \emph{Stochastic RAG: End-to-End Retrieval-Augmented Generation through Expected Utility Maximization}.
\newblock arXiv:2405.02816v1, 2024.
\newblock \url{https://arxiv.org/abs/2405.02816v1}

\bibitem{arxiv:2208.03299v3}
Gautier Izacard, Patrick Lewis, Maria Lomeli, Lucas Hosseini, Fabio Petroni, Timo Schick \emph{et al.}.
\newblock \emph{Atlas Few-shot Learning with Retrieval Augmented Language Models}.
\newblock arXiv:2208.03299v3, 2022.
\newblock \url{https://arxiv.org/abs/2208.03299v3}

\bibitem{arxiv:2310.01558v1}
Ori Yoran, Tomer Wolfson, Ori Ram, Jonathan Berant.
\newblock \emph{Making Retrieval-Augmented Language Models Robust to Irrelevant Context}.
\newblock arXiv:2310.01558v1, 2023.
\newblock \url{https://arxiv.org/abs/2310.01558v1}

\bibitem{arxiv:2401.05856v1}
Scott Barnett, Stefanus Kurniawan, Srikanth Thudumu, Zach Brannelly, Mohamed Abdelrazek.
\newblock \emph{Seven Failure Points When Engineering a Retrieval Augmented Generation System}.
\newblock arXiv:2401.05856v1, 2024.
\newblock \url{https://arxiv.org/abs/2401.05856v1}

\bibitem{arxiv:2305.06983v2}
Zhengbao Jiang, Frank F. Xu, Luyu Gao, Zhiqing Sun, Qian Liu, Jane Dwivedi-Yu \emph{et al.}.
\newblock \emph{Active Retrieval Augmented Generation}.
\newblock arXiv:2305.06983v2, 2023.
\newblock \url{https://arxiv.org/abs/2305.06983v2}

\bibitem{arxiv:2405.13792v1}
Xin Cheng, Xun Wang, Xingxing Zhang, Tao Ge, Si-Qing Chen, Furu Wei \emph{et al.}.
\newblock \emph{xRAG: Extreme Context Compression for Retrieval-augmented Generation with One Token}.
\newblock arXiv:2405.13792v1, 2024.
\newblock \url{https://arxiv.org/abs/2405.13792v1}

\bibitem{arxiv:2408.08921v2}
Boci Peng, Yun Zhu, Yongchao Liu, Xiaohe Bo, Haizhou Shi, Chuntao Hong \emph{et al.}.
\newblock \emph{Graph Retrieval-Augmented Generation: A Survey}.
\newblock arXiv:2408.08921v2, 2024.
\newblock \url{https://arxiv.org/abs/2408.08921v2}

\bibitem{arxiv:2210.02928v2}
Wenhu Chen, Hexiang Hu, Xi Chen, Pat Verga, William W. Cohen.
\newblock \emph{MuRAG Multimodal Retrieval-Augmented Generator for Open Question Answering over Images and Text}.
\newblock arXiv:2210.02928v2, 2022.
\newblock \url{https://arxiv.org/abs/2210.02928v2}

\bibitem{arxiv:2211.12561v2}
Michihiro Yasunaga, Armen Aghajanyan, Weijia Shi, Rich James, Jure Leskovec, Percy Liang \emph{et al.}.
\newblock \emph{Retrieval-Augmented Multimodal Language Modeling}.
\newblock arXiv:2211.12561v2, 2022.
\newblock \url{https://arxiv.org/abs/2211.12561v2}

\bibitem{arxiv:2401.13256v1}
Hongru Wang, Wenyu Huang, Yang Deng, Rui Wang, Zezhong Wang, Yufei Wang \emph{et al.}.
\newblock \emph{UniMS-RAG A Unified Multi-source Retrieval-Augmented Generation for Personalized Dialogue Systems}.
\newblock arXiv:2401.13256v1, 2024.
\newblock \url{https://arxiv.org/abs/2401.13256v1}

\bibitem{arxiv:2305.15294v2}
Zhihong Shao, Yeyun Gong, Yelong Shen, Minlie Huang, Nan Duan, Weizhu Chen.
\newblock \emph{Enhancing Retrieval-Augmented Large Language Models with Iterative Retrieval-Generation Synergy}.
\newblock arXiv:2305.15294v2, 2023.
\newblock \url{https://arxiv.org/abs/2305.15294v2}

\bibitem{arxiv:2402.11626v1}
Yujia Zhou, Zheng Liu, Jiajie Jin, Jian-Yun Nie, Zhicheng Dou.
\newblock \emph{Metacognitive Retrieval-Augmented Large Language Models}.
\newblock arXiv:2402.11626v1, 2024.
\newblock \url{https://arxiv.org/abs/2402.11626v1}

\bibitem{arxiv:2212.10509v2}
Harsh Trivedi, Niranjan Balasubramanian, Tushar Khot, Ashish Sabharwal.
\newblock \emph{Interleaving Retrieval with Chain-of-Thought Reasoning for Knowledge-Intensive Multi-Step Questions}.
\newblock arXiv:2212.10509v2, 2022.
\newblock \url{https://arxiv.org/abs/2212.10509v2}

\bibitem{arxiv:2401.15884v2}
Shi-Qi Yan, Jia-Chen Gu, Yun Zhu, Zhen-Hua Ling.
\newblock \emph{Corrective Retrieval Augmented Generation}.
\newblock arXiv:2401.15884v2, 2024.
\newblock \url{https://arxiv.org/abs/2401.15884v2}

\bibitem{arxiv:2402.13492v3}
Seiji Maekawa, Hayate Iso, Sairam Gurajada, Nikita Bhutani.
\newblock \emph{Retrieval Helps or Hurts A Deeper Dive into the Efficacy of Retrieval Augmentation to Language Models}.
\newblock arXiv:2402.13492v3, 2024.
\newblock \url{https://arxiv.org/abs/2402.13492v3}

\bibitem{arxiv:2209.14290v1}
Sebastian Hofst\"{a}tter, Jiecao Chen, Karthik Raman, Hamed Zamani.
\newblock \emph{FiD-Light Efficient and Effective Retrieval-Augmented Text Generation}.
\newblock arXiv:2209.14290v1, 2022.
\newblock \url{https://arxiv.org/abs/2209.14290v1}

\bibitem{arxiv:2402.01176v2}
Xiaoxi Li, Zhicheng Dou, Yujia Zhou, Fangchao Liu.
\newblock \emph{CorpusLM Towards a Unified Language Model on Corpus for Knowledge-Intensive Tasks}.
\newblock arXiv:2402.01176v2, 2024.
\newblock \url{https://arxiv.org/abs/2402.01176v2}

\bibitem{arxiv:2207.06300v1}
Michael Glass, Gaetano Rossiello, Md Faisal Mahbub Chowdhury, Ankita Rajaram Naik, Pengshan Cai, Alfio Gliozzo.
\newblock \emph{Re2G Retrieve, Rerank, Generate}.
\newblock arXiv:2207.06300v1, 2022.
\newblock \url{https://arxiv.org/abs/2207.06300v1}

\bibitem{arxiv:2312.10997v5}
Yunfan Gao, Yun Xiong, Xinyu Gao, Kangxiang Jia, Jinliu Pan, Yuxi Bi \emph{et al.}.
\newblock \emph{Retrieval-Augmented Generation for Large Language Models A Survey}.
\newblock arXiv:2312.10997v5, 2023.
\newblock \url{https://arxiv.org/abs/2312.10997v5}

\bibitem{arxiv:2406.15319v3}
Ziyan Jiang, Xueguang Ma, Wenhu Chen.
\newblock \emph{LongRAG: Enhancing Retrieval-Augmented Generation with Long-context LLMs}.
\newblock arXiv:2406.15319v3, 2024.
\newblock \url{https://arxiv.org/abs/2406.15319v3}

\bibitem{arxiv:2405.06683v1}
Yunxiao Shi, Xing Zi, Zijing Shi, Haimin Zhang, Qiang Wu, Min Xu.
\newblock \emph{ERAGent: Enhancing Retrieval-Augmented Language Models with Improved Accuracy, Efficiency, and Personalization}.
\newblock arXiv:2405.06683v1, 2024.
\newblock \url{https://arxiv.org/abs/2405.06683v1}

\bibitem{arxiv:2409.05591v2}
Hongjin Qian, Peitian Zhang, Zheng Liu, Kelong Mao, Zhicheng Dou.
\newblock \emph{MemoRAG: Moving towards Next-Gen RAG Via Memory-Inspired Knowledge Discovery}.
\newblock arXiv:2409.05591v2, 2024.
\newblock \url{https://arxiv.org/abs/2409.05591v2}

\bibitem{arxiv:2404.14851v3}
Xiaoxi Li, Jiajie Jin, Yujia Zhou, Yuyao Zhang, Peitian Zhang, Yutao Zhu \emph{et al.}.
\newblock \emph{From Matching to Generation: A Survey on Generative Information Retrieval}.
\newblock arXiv:2404.14851v3, 2024.
\newblock \url{https://arxiv.org/abs/2404.14851v3}

\bibitem{arxiv:2401.06532v2}
Yutao Zhu, Peitian Zhang, Chenghao Zhang, Yifei Chen, Binyu Xie, Zhicheng Dou \emph{et al.}.
\newblock \emph{INTERS Unlocking the Power of Large Language Models in Search with Instruction Tuning}.
\newblock arXiv:2401.06532v2, 2024.
\newblock \url{https://arxiv.org/abs/2401.06532v2}

\bibitem{arxiv:2403.15246v1}
Orion Weller, Benjamin Chang, Sean MacAvaney, Kyle Lo, Arman Cohan, Benjamin Van Durme \emph{et al.}.
\newblock \emph{FollowIR Evaluating and Teaching Information Retrieval Models to Follow Instructions}.
\newblock arXiv:2403.15246v1, 2024.
\newblock \url{https://arxiv.org/abs/2403.15246v1}

\bibitem{arxiv:2403.05676v1}
Wenqi Jiang, Shuai Zhang, Boran Han, Jie Wang, Bernie Wang, Tim Kraska.
\newblock \emph{PipeRAG Fast Retrieval-Augmented Generation via Algorithm-System Co-design}.
\newblock arXiv:2403.05676v1, 2024.
\newblock \url{https://arxiv.org/abs/2403.05676v1}

\bibitem{arxiv:2401.14021v1}
Zhihao Zhang, Alan Zhu, Lijie Yang, Yihua Xu, Lanting Li, Phitchaya Mangpo Phothilimthana \emph{et al.}.
\newblock \emph{Accelerating Retrieval-Augmented Language Model Serving with Speculation}.
\newblock arXiv:2401.14021v1, 2024.
\newblock \url{https://arxiv.org/abs/2401.14021v1}

\bibitem{arxiv:2405.13576v1}
Jiajie Jin, Yutao Zhu, Xinyu Yang, Chenghao Zhang, Zhicheng Dou.
\newblock \emph{FlashRAG: A Modular Toolkit for Efficient Retrieval-Augmented Generation Research}.
\newblock arXiv:2405.13576v1, 2024.
\newblock \url{https://arxiv.org/abs/2405.13576v1}

\bibitem{arxiv:2402.16767v1}
Jiafeng Guo, Changjiang Zhou, Ruqing Zhang, Jiangui Chen, Maarten de Rijke, Yixing Fan \emph{et al.}.
\newblock \emph{CorpusBrain++ A Continual Generative Pre-Training Framework for Knowledge-Intensive Language Tasks}.
\newblock arXiv:2402.16767v1, 2024.
\newblock \url{https://arxiv.org/abs/2402.16767v1}

\bibitem{arxiv:2404.14600v1}
Hansi Zeng, Chen Luo, Hamed Zamani.
\newblock \emph{Planning Ahead in Generative Retrieval Guiding Autoregressive Generation through Simultaneous Decoding}.
\newblock arXiv:2404.14600v1, 2024.
\newblock \url{https://arxiv.org/abs/2404.14600v1}

\bibitem{arxiv:2009.08553v4}
Yuning Mao, Pengcheng He, Xiaodong Liu, Yelong Shen, Jianfeng Gao, Jiawei Han \emph{et al.}.
\newblock \emph{Generation-Augmented Retrieval for Open-domain Question Answering}.
\newblock arXiv:2009.08553v4, 2020.
\newblock \url{https://arxiv.org/abs/2009.08553v4}

\bibitem{arxiv:2408.01875v2}
Yanfei Chen, Jinsung Yoon, Devendra Singh Sachan, Qingze Wang, Vincent Cohen-Addad, Mohammadhossein Bateni \emph{et al.}.
\newblock \emph{Re-Invoke: Tool Invocation Rewriting for Zero-Shot Tool Retrieval}.
\newblock arXiv:2408.01875v2, 2024.
\newblock \url{https://arxiv.org/abs/2408.01875v2}

\bibitem{arxiv:2404.04302v1}
Nirmalie Wiratunga, Ramitha Abeyratne, Lasal Jayawardena, Kyle Martin, Stewart Massie, Ikechukwu Nkisi-Orji \emph{et al.}.
\newblock \emph{CBR-RAG Case-Based Reasoning for Retrieval Augmented Generation in LLMs for Legal Question Answering}.
\newblock arXiv:2404.04302v1, 2024.
\newblock \url{https://arxiv.org/abs/2404.04302v1}

\bibitem{arxiv:2406.14282v1}
Junjie Wang, Mingyang Chen, Binbin Hu, Dan Yang, Ziqi Liu, Yue Shen \emph{et al.}.
\newblock \emph{Learning to Plan for Retrieval-Augmented Large Language Models from Knowledge Graphs}.
\newblock arXiv:2406.14282v1, 2024.
\newblock \url{https://arxiv.org/abs/2406.14282v1}

\bibitem{arxiv:2404.13556v1}
Kelong Mao, Chenlong Deng, Haonan Chen, Fengran Mo, Zheng Liu, Tetsuya Sakai \emph{et al.}.
\newblock \emph{ChatRetriever Adapting Large Language Models for Generalized and Robust Conversational Dense Retrieval}.
\newblock arXiv:2404.13556v1, 2024.
\newblock \url{https://arxiv.org/abs/2404.13556v1}

\bibitem{arxiv:2209.11755v1}
Zhuyun Dai, Vincent Y. Zhao, Ji Ma, Yi Luan, Jianmo Ni, Jing Lu \emph{et al.}.
\newblock \emph{Promptagator Few-shot Dense Retrieval From 8 Examples}.
\newblock arXiv:2209.11755v1, 2022.
\newblock \url{https://arxiv.org/abs/2209.11755v1}

\bibitem{arxiv:2401.13222v2}
Anoushka Gade, Jorjeta Jetcheva.
\newblock \emph{It's About Time Incorporating Temporality in Retrieval Augmented Language Models}.
\newblock arXiv:2401.13222v2, 2024.
\newblock \url{https://arxiv.org/abs/2401.13222v2}

\bibitem{arxiv:2309.01431v2}
Jiawei Chen, Hongyu Lin, Xianpei Han, Le Sun.
\newblock \emph{Benchmarking Large Language Models in Retrieval-Augmented Generation}.
\newblock arXiv:2309.01431v2, 2023.
\newblock \url{https://arxiv.org/abs/2309.01431v2}

\bibitem{arxiv:2401.14887v3}
Florin Cuconasu, Giovanni Trappolini, Federico Siciliano, Simone Filice, Cesare Campagnano, Yoelle Maarek \emph{et al.}.
\newblock \emph{The Power of Noise Redefining Retrieval for RAG Systems}.
\newblock arXiv:2401.14887v3, 2024.
\newblock \url{https://arxiv.org/abs/2401.14887v3}

\bibitem{arxiv:2402.11129v1}
Haoyu Wang, Tuo Zhao, Jing Gao.
\newblock \emph{BlendFilter Advancing Retrieval-Augmented Large Language Models via Query Generation Blending and Knowledge Filtering}.
\newblock arXiv:2402.11129v1, 2024.
\newblock \url{https://arxiv.org/abs/2402.11129v1}

\bibitem{arxiv:2402.16457v1}
Zihan Zhang, Meng Fang, Ling Chen.
\newblock \emph{RetrievalQA Assessing Adaptive Retrieval-Augmented Generation for Short-form Open-Domain Question Answering}.
\newblock arXiv:2402.16457v1, 2024.
\newblock \url{https://arxiv.org/abs/2402.16457v1}

\bibitem{arxiv:2403.06840v1}
Yanming Liu, Xinyue Peng, Xuhong Zhang, Weihao Liu, Jianwei Yin, Jiannan Cao \emph{et al.}.
\newblock \emph{RA-ISF Learning to Answer and Understand from Retrieval Augmentation via Iterative Self-Feedback}.
\newblock arXiv:2403.06840v1, 2024.
\newblock \url{https://arxiv.org/abs/2403.06840v1}

\bibitem{arxiv:2402.12317v1}
Hongjin Su, Shuyang Jiang, Yuhang Lai, Haoyuan Wu, Boao Shi, Che Liu \emph{et al.}.
\newblock \emph{ARKS Active Retrieval in Knowledge Soup for Code Generation}.
\newblock arXiv:2402.12317v1, 2024.
\newblock \url{https://arxiv.org/abs/2402.12317v1}

\bibitem{arxiv:2306.05212v1}
Jiongnan Liu, Jiajie Jin, Zihan Wang, Jiehan Cheng, Zhicheng Dou, Ji-Rong Wen.
\newblock \emph{RETA-LLM A Retrieval-Augmented Large Language Model Toolkit}.
\newblock arXiv:2306.05212v1, 2023.
\newblock \url{https://arxiv.org/abs/2306.05212v1}

\bibitem{arxiv:2401.15391v1}
Yixuan Tang, Yi Yang.
\newblock \emph{MultiHop-RAG Benchmarking Retrieval-Augmented Generation for Multi-Hop Queries}.
\newblock arXiv:2401.15391v1, 2024.
\newblock \url{https://arxiv.org/abs/2401.15391v1}

\bibitem{arxiv:2309.15217v1}
Shahul Es, Jithin James, Luis Espinosa-Anke, Steven Schockaert.
\newblock \emph{RAGAS Automated Evaluation of Retrieval Augmented Generation}.
\newblock arXiv:2309.15217v1, 2023.
\newblock \url{https://arxiv.org/abs/2309.15217v1}

\bibitem{arxiv:2305.17331v1}
Zichun Yu, Chenyan Xiong, Shi Yu, Zhiyuan Liu.
\newblock \emph{Augmentation-Adapted Retriever Improves Generalization of Language Models as Generic Plug-In}.
\newblock arXiv:2305.17331v1, 2023.
\newblock \url{https://arxiv.org/abs/2305.17331v1}

\bibitem{arxiv:2402.01733v1}
YuHe Ke, Liyuan Jin, Kabilan Elangovan, Hairil Rizal Abdullah, Nan Liu, Alex Tiong Heng Sia \emph{et al.}.
\newblock \emph{Development and Testing of Retrieval Augmented Generation in Large Language Models -- A Case Study Report}.
\newblock arXiv:2402.01733v1, 2024.
\newblock \url{https://arxiv.org/abs/2402.01733v1}

\bibitem{arxiv:2402.18150v1}
Shicheng Xu, Liang Pang, Mo Yu, Fandong Meng, Huawei Shen, Xueqi Cheng \emph{et al.}.
\newblock \emph{Unsupervised Information Refinement Training of Large Language Models for Retrieval-Augmented Generation}.
\newblock arXiv:2402.18150v1, 2024.
\newblock \url{https://arxiv.org/abs/2402.18150v1}

\bibitem{arxiv:2406.14891v2}
Zhengliang Shi, Weiwei Sun, Shen Gao, Pengjie Ren, Zhumin Chen, Zhaochun Ren.
\newblock \emph{Generate-then-Ground in Retrieval-Augmented Generation for Multi-hop Question Answering}.
\newblock arXiv:2406.14891v2, 2024.
\newblock \url{https://arxiv.org/abs/2406.14891v2}

\bibitem{arxiv:2110.07752v2}
Ashwin Paranjape, Omar Khattab, Christopher Potts, Matei Zaharia, Christopher D. Manning.
\newblock \emph{Hindsight Posterior-guided training of retrievers for improved open-ended generation}.
\newblock arXiv:2110.07752v2, 2021.
\newblock \url{https://arxiv.org/abs/2110.07752v2}

\bibitem{arxiv:2402.12174v1}
Jiajie Jin, Yutao Zhu, Yujia Zhou, Zhicheng Dou.
\newblock \emph{BIDER Bridging Knowledge Inconsistency for Efficient Retrieval-Augmented LLMs via Key Supporting Evidence}.
\newblock arXiv:2402.12174v1, 2024.
\newblock \url{https://arxiv.org/abs/2402.12174v1}

\bibitem{arxiv:2404.10981v1}
Yizheng Huang, Jimmy Huang.
\newblock \emph{A Survey on Retrieval-Augmented Text Generation for Large Language Models}.
\newblock arXiv:2404.10981v1, 2024.
\newblock \url{https://arxiv.org/abs/2404.10981v1}

\bibitem{arxiv:2405.07437v2}
Hao Yu, Aoran Gan, Kai Zhang, Shiwei Tong, Qi Liu, Zhaofeng Liu.
\newblock \emph{Evaluation of Retrieval-Augmented Generation: A Survey}.
\newblock arXiv:2405.07437v2, 2024.
\newblock \url{https://arxiv.org/abs/2405.07437v2}

\bibitem{arxiv:2404.13781v1}
Alireza Salemi, Hamed Zamani.
\newblock \emph{Evaluating Retrieval Quality in Retrieval-Augmented Generation}.
\newblock arXiv:2404.13781v1, 2024.
\newblock \url{https://arxiv.org/abs/2404.13781v1}

\bibitem{arxiv:2307.11019v2}
Ruiyang Ren, Yuhao Wang, Yingqi Qu, Wayne Xin Zhao, Jing Liu, Hao Tian \emph{et al.}.
\newblock \emph{Investigating the Factual Knowledge Boundary of Large Language Models with Retrieval Augmentation}.
\newblock arXiv:2307.11019v2, 2023.
\newblock \url{https://arxiv.org/abs/2307.11019v2}

\bibitem{arxiv:2406.14497v1}
Zora Zhiruo Wang, Akari Asai, Xinyan Velocity Yu, Frank F. Xu, Yiqing Xie, Graham Neubig \emph{et al.}.
\newblock \emph{CodeRAG-Bench: Can Retrieval Augment Code Generation?}.
\newblock arXiv:2406.14497v1, 2024.
\newblock \url{https://arxiv.org/abs/2406.14497v1}

\bibitem{arxiv:2311.09476v2}
Jon Saad-Falcon, Omar Khattab, Christopher Potts, Matei Zaharia.
\newblock \emph{ARES An Automated Evaluation Framework for Retrieval-Augmented Generation Systems}.
\newblock arXiv:2311.09476v2, 2023.
\newblock \url{https://arxiv.org/abs/2311.09476v2}

\bibitem{arxiv:2405.20446v2}
Maya Anderson, Guy Amit, Abigail Goldsteen.
\newblock \emph{Is My Data in Your Retrieval Database? Membership Inference Attacks Against Retrieval Augmented Generation}.
\newblock arXiv:2405.20446v2, 2024.
\newblock \url{https://arxiv.org/abs/2405.20446v2}

\bibitem{arxiv:2405.20978v1}
Feiteng Fang, Yuelin Bai, Shiwen Ni, Min Yang, Xiaojun Chen, Ruifeng Xu.
\newblock \emph{Enhancing Noise Robustness of Retrieval-Augmented Language Models with Adaptive Adversarial Training}.
\newblock arXiv:2405.20978v1, 2024.
\newblock \url{https://arxiv.org/abs/2405.20978v1}

\bibitem{arxiv:2407.08223v1}
Zilong Wang, Zifeng Wang, Long Le, Huaixiu Steven Zheng, Swaroop Mishra, Vincent Perot \emph{et al.}.
\newblock \emph{Speculative RAG: Enhancing Retrieval Augmented Generation through Drafting}.
\newblock arXiv:2407.08223v1, 2024.
\newblock \url{https://arxiv.org/abs/2407.08223v1}

\bibitem{arxiv:2302.12813v3}
Baolin Peng, Michel Galley, Pengcheng He, Hao Cheng, Yujia Xie, Yu Hu \emph{et al.}.
\newblock \emph{Check Your Facts and Try Again Improving Large Language Models with External Knowledge and Automated Feedback}.
\newblock arXiv:2302.12813v3, 2023.
\newblock \url{https://arxiv.org/abs/2302.12813v3}

\bibitem{arxiv:2402.10612v1}
Hanxing Ding, Liang Pang, Zihao Wei, Huawei Shen, Xueqi Cheng.
\newblock \emph{Retrieve Only When It Needs Adaptive Retrieval Augmentation for Hallucination Mitigation in Large Language Models}.
\newblock arXiv:2402.10612v1, 2024.
\newblock \url{https://arxiv.org/abs/2402.10612v1}

\bibitem{arxiv:1809.05296v5}
Deng Cai, Yan Wang, Victoria Bi, Zhaopeng Tu, Xiaojiang Liu, Wai Lam \emph{et al.}.
\newblock \emph{Skeleton-to-Response Dialogue Generation Guided by Retrieval Memory}.
\newblock arXiv:1809.05296v5, 2018.
\newblock \url{https://arxiv.org/abs/1809.05296v5}

\bibitem{arxiv:2405.13021v1}
Diji Yang, Jinmeng Rao, Kezhen Chen, Xiaoyuan Guo, Yawen Zhang, Jie Yang \emph{et al.}.
\newblock \emph{IM-RAG: Multi-Round Retrieval-Augmented Generation Through Learning Inner Monologues}.
\newblock arXiv:2405.13021v1, 2024.
\newblock \url{https://arxiv.org/abs/2405.13021v1}

\bibitem{arxiv:2408.04948v1}
Bhaskarjit Sarmah, Benika Hall, Rohan Rao, Sunil Patel, Stefano Pasquali, Dhagash Mehta.
\newblock \emph{HybridRAG: Integrating Knowledge Graphs and Vector Retrieval Augmented Generation for Efficient Information Extraction}.
\newblock arXiv:2408.04948v1, 2024.
\newblock \url{https://arxiv.org/abs/2408.04948v1}

\bibitem{arxiv:2108.11601v2}
Md Rizwan Parvez, Wasi Uddin Ahmad, Saikat Chakraborty, Baishakhi Ray, Kai-Wei Chang.
\newblock \emph{Retrieval Augmented Code Generation and Summarization}.
\newblock arXiv:2108.11601v2, 2021.
\newblock \url{https://arxiv.org/abs/2108.11601v2}

\bibitem{arxiv:2407.12883v1}
Hongjin Su, Howard Yen, Mengzhou Xia, Weijia Shi, Niklas Muennighoff, Han-yu Wang \emph{et al.}.
\newblock \emph{BRIGHT: A Realistic and Challenging Benchmark for Reasoning-Intensive Retrieval}.
\newblock arXiv:2407.12883v1, 2024.
\newblock \url{https://arxiv.org/abs/2407.12883v1}

\bibitem{arxiv:2409.10516v2}
Di Liu, Meng Chen, Baotong Lu, Huiqiang Jiang, Zhenhua Han, Qianxi Zhang \emph{et al.}.
\newblock \emph{RetrievalAttention: Accelerating Long-Context LLM Inference via Vector Retrieval}.
\newblock arXiv:2409.10516v2, 2024.
\newblock \url{https://arxiv.org/abs/2409.10516v2}

\bibitem{arxiv:2402.13178v2}
Guangzhi Xiong, Qiao Jin, Zhiyong Lu, Aidong Zhang.
\newblock \emph{Benchmarking Retrieval-Augmented Generation for Medicine}.
\newblock arXiv:2402.13178v2, 2024.
\newblock \url{https://arxiv.org/abs/2402.13178v2}

\bibitem{arxiv:2404.07220v1}
Kunal Sawarkar, Abhilasha Mangal, Shivam Raj Solanki.
\newblock \emph{Blended RAG Improving RAG (Retriever-Augmented Generation) Accuracy with Semantic Search and Hybrid Query-Based Retrievers}.
\newblock arXiv:2404.07220v1, 2024.
\newblock \url{https://arxiv.org/abs/2404.07220v1}

\bibitem{arxiv:2406.05085v1}
Maciej Besta, Ales Kubicek, Roman Niggli, Robert Gerstenberger, Lucas Weitzendorf, Mingyuan Chi \emph{et al.}.
\newblock \emph{Multi-Head RAG: Solving Multi-Aspect Problems with LLMs}.
\newblock arXiv:2406.05085v1, 2024.
\newblock \url{https://arxiv.org/abs/2406.05085v1}

\bibitem{arxiv:2407.21439v2}
Zhanpeng Chen, Chengjin Xu, Yiyan Qi, Jian Guo.
\newblock \emph{MLLM Is a Strong Reranker: Advancing Multimodal Retrieval-augmented Generation via Knowledge-enhanced Reranking and Noise-injected Training}.
\newblock arXiv:2407.21439v2, 2024.
\newblock \url{https://arxiv.org/abs/2407.21439v2}

\bibitem{arxiv:2404.17723v2}
Zhentao Xu, Mark Jerome Cruz, Matthew Guevara, Tie Wang, Manasi Deshpande, Xiaofeng Wang \emph{et al.}.
\newblock \emph{Retrieval-Augmented Generation with Knowledge Graphs for Customer Service Question Answering}.
\newblock arXiv:2404.17723v2, 2024.
\newblock \url{https://arxiv.org/abs/2404.17723v2}

\bibitem{arxiv:2404.18443v1}
Ran Xu, Wenqi Shi, Yue Yu, Yuchen Zhuang, Yanqiao Zhu, May D. Wang \emph{et al.}.
\newblock \emph{BMRetriever: Tuning Large Language Models as Better Biomedical Text Retrievers}.
\newblock arXiv:2404.18443v1, 2024.
\newblock \url{https://arxiv.org/abs/2404.18443v1}

\bibitem{arxiv:2402.04853v1}
Ekaterina Khramtsova, Shengyao Zhuang, Mahsa Baktashmotlagh, Guido Zuccon.
\newblock \emph{Leveraging LLMs for Unsupervised Dense Retriever Ranking}.
\newblock arXiv:2402.04853v1, 2024.
\newblock \url{https://arxiv.org/abs/2402.04853v1}

\bibitem{arxiv:2305.02437v3}
Xin Cheng, Di Luo, Xiuying Chen, Lemao Liu, Dongyan Zhao, Rui Yan.
\newblock \emph{Lift Yourself Up Retrieval-augmented Text Generation with Self Memory}.
\newblock arXiv:2305.02437v3, 2023.
\newblock \url{https://arxiv.org/abs/2305.02437v3}

\bibitem{arxiv:2403.14403v2}
Soyeong Jeong, Jinheon Baek, Sukmin Cho, Sung Ju Hwang, Jong C. Park.
\newblock \emph{Adaptive-RAG Learning to Adapt Retrieval-Augmented Large Language Models through Question Complexity}.
\newblock arXiv:2403.14403v2, 2024.
\newblock \url{https://arxiv.org/abs/2403.14403v2}

\bibitem{arxiv:2407.21059v1}
Yunfan Gao, Yun Xiong, Meng Wang, Haofen Wang.
\newblock \emph{Modular RAG: Transforming RAG Systems into LEGO-like Reconfigurable Frameworks}.
\newblock arXiv:2407.21059v1, 2024.
\newblock \url{https://arxiv.org/abs/2407.21059v1}

\bibitem{arxiv:2404.11672v1}
Ali Modarressi, Abdullatif K\"{o}ksal, Ayyoob Imani, Mohsen Fayyaz, Hinrich Sch\"{u}tze.
\newblock \emph{MemLLM Finetuning LLMs to Use An Explicit Read-Write Memory}.
\newblock arXiv:2404.11672v1, 2024.
\newblock \url{https://arxiv.org/abs/2404.11672v1}

\bibitem{arxiv:2112.09737v2}
Niket Tandon, Aman Madaan, Peter Clark, Yiming Yang.
\newblock \emph{Learning to Repair Repairing model output errors after deployment using a dynamic memory of feedback}.
\newblock arXiv:2112.09737v2, 2021.
\newblock \url{https://arxiv.org/abs/2112.09737v2}

\bibitem{arxiv:2206.06520v1}
Eric Mitchell, Charles Lin, Antoine Bosselut, Christopher D. Manning, Chelsea Finn.
\newblock \emph{Memory-Based Model Editing at Scale}.
\newblock arXiv:2206.06520v1, 2022.
\newblock \url{https://arxiv.org/abs/2206.06520v1}

\bibitem{arxiv:2205.09707v1}
Keshav Santhanam, Omar Khattab, Christopher Potts, Matei Zaharia.
\newblock \emph{PLAID An Efficient Engine for Late Interaction Retrieval}.
\newblock arXiv:2205.09707v1, 2022.
\newblock \url{https://arxiv.org/abs/2205.09707v1}

\bibitem{arxiv:2206.02873v5}
Guilherme Moraes Rosa, Luiz Bonifacio, Vitor Jeronymo, Hugo Abonizio, Marzieh Fadaee, Roberto Lotufo \emph{et al.}.
\newblock \emph{No Parameter Left Behind How Distillation and Model Size Affect Zero-Shot Retrieval}.
\newblock arXiv:2206.02873v5, 2022.
\newblock \url{https://arxiv.org/abs/2206.02873v5}

\bibitem{arxiv:2208.07652v1}
Jiangui Chen, Ruqing Zhang, Jiafeng Guo, Yiqun Liu, Yixing Fan, Xueqi Cheng.
\newblock \emph{CorpusBrain Pre-train a Generative Retrieval Model for Knowledge-Intensive Language Tasks}.
\newblock arXiv:2208.07652v1, 2022.
\newblock \url{https://arxiv.org/abs/2208.07652v1}

\bibitem{arxiv:2305.11841v1}
Ronak Pradeep, Kai Hui, Jai Gupta, Adam D. Lelkes, Honglei Zhuang, Jimmy Lin \emph{et al.}.
\newblock \emph{How Does Generative Retrieval Scale to Millions of Passages}.
\newblock arXiv:2305.11841v1, 2023.
\newblock \url{https://arxiv.org/abs/2305.11841v1}

\bibitem{arxiv:2405.07767v1}
Hossein A. Rahmani, Nick Craswell, Emine Yilmaz, Bhaskar Mitra, Daniel Campos.
\newblock \emph{Synthetic Test Collections for Retrieval Evaluation}.
\newblock arXiv:2405.07767v1, 2024.
\newblock \url{https://arxiv.org/abs/2405.07767v1}

\bibitem{arxiv:2407.16833v1}
Zhuowan Li, Cheng Li, Mingyang Zhang, Qiaozhu Mei, Michael Bendersky.
\newblock \emph{Retrieval Augmented Generation or Long-Context LLMs? A Comprehensive Study and Hybrid Approach}.
\newblock arXiv:2407.16833v1, 2024.
\newblock \url{https://arxiv.org/abs/2407.16833v1}

\end{thebibliography}

\end{document}
